\documentclass[aspectratio=169,10pt]{beamer}
% notes.tex defines \NOTESONLY and reads this file: then only the speaker notes are printed
\ifdefined\NOTESONLY
  \setbeameroption{show only notes}
\fi

\usepackage{lmodern}
\usepackage[utf8]{inputenc}
\usepackage[T1]{fontenc}
\usepackage[polish,english]{babel}
\usepackage{amsmath,amssymb}
\usepackage{booktabs}
\usepackage{colortbl}
\usepackage{xcolor}
\usepackage{fontawesome5}
\usepackage{listings}
\usepackage{etoolbox}
\usepackage{tikz}
\usetikzlibrary{arrows.meta, calc, positioning, fit, decorations.pathreplacing, decorations.pathmorphing}
\usepackage{pgfplots}
% layers for the code windows drawn in TikZ: the window, then the line that runs now, then the text
\pgfdeclarelayer{codebg}\pgfdeclarelayer{codehl}\pgfsetlayers{codebg,codehl,main}
\pgfplotsset{compat=1.18}

\usetheme{Madrid}
\useinnertheme{rounded}
\useoutertheme{miniframes}
\setbeamertemplate{navigation symbols}{}
\setbeamertemplate{itemize items}[circle]
\setbeamertemplate{enumerate items}[default]
\setbeamertemplate{blocks}[rounded][shadow=true]
% one dot per idea, not per frame: \samedot before a frame gives it the dot of the frame before
% (it undoes the step of the dot counter that the next frame makes)
\newcommand{\samedot}{\addtocounter{subsectionslide}{-1}}
% the default dot step of the circle dots (0.1 cm dots)
\setbeamersize{mini frame size=.14cm,mini frame offset=-.01cm}

% ======================================================================
% Colour rules for presentations (Beamer + TikZ), from Agents/colors
% ======================================================================

% ---------- template ----------
\definecolor{DeepBlue}{HTML}{283593}      % title bar, headings, links, control-flow arrows
\definecolor{AccentOrange}{HTML}{EB5A14}  % annotations: step numbers, "look here", code highlight, conclusions
\definecolor{LightGray}{HTML}{F5F6FA}     % background of information blocks
\setbeamercolor{palette primary}{bg=DeepBlue,fg=white}
\setbeamercolor{palette secondary}{bg=DeepBlue!80!black,fg=white}
\setbeamercolor{palette tertiary}{bg=DeepBlue!90!black,fg=white}
\setbeamercolor{title}{bg=DeepBlue,fg=white}
\setbeamercolor{frametitle}{bg=DeepBlue,fg=white}
\setbeamercolor{structure}{fg=DeepBlue}
\setbeamercolor{item}{fg=AccentOrange}
\setbeamercolor{alerted text}{fg=AccentOrange}
\setbeamercolor{block title}{bg=DeepBlue,fg=white}
\setbeamercolor{block body}{bg=LightGray,fg=black}
\setbeamercolor{block title example}{bg=AccentOrange,fg=white}
\setbeamercolor{block body example}{bg=AccentOrange!10,fg=black}
\setbeamercolor{enumerate item}{fg=DeepBlue}
\setbeamerfont{enumerate item}{series=\bfseries}

% ---------- actors and categories: fills only (S1 to S6) ----------
\definecolor{sIndigo}{HTML}{5C6BC0}   % S1
\definecolor{sPink}{HTML}{EC407A}     % S2
\definecolor{sViolet}{HTML}{7E57C2}   % S3
\definecolor{sAmber}{HTML}{FFB300}    % S4: black text on it
\definecolor{sTeal}{HTML}{26A69A}     % S5
\definecolor{sBrown}{HTML}{8D6E63}    % S6
% text variants that stay readable on white
\definecolor{sPinkText}{HTML}{C2185B}
\definecolor{sAmberText}{HTML}{A06A00}
\definecolor{sTealText}{HTML}{00796B}

% ---------- actions: frames only; results: signs only ----------
\definecolor{aExec}{HTML}{43A047}     % runs now
\definecolor{aRead}{HTML}{1E88E5}     % is read
\definecolor{aWrite}{HTML}{E53935}    % is written

% ---------- neutral (Blue Grey) ----------
\definecolor{nInk}{HTML}{37474F}      % text in figures, waveforms, wires, data-flow arrows
\definecolor{nLine}{HTML}{90A4AE}     % borders, axes
\definecolor{nLight}{HTML}{CFD8DC}    % inactive lines, grids, empty places (dashed)
\definecolor{nFill}{HTML}{ECEFF1}     % fill of boxes and used cells
\definecolor{nIdle}{HTML}{E0E0E0}     % time lost (idle, busy waiting, bubbles)

% ---------- TikZ styles ----------
\tikzset{
  aexec/.style={draw=aExec, line width=1.6pt},
  aexecm/.style={aexec, label={[text=aExec, inner sep=1pt, font=\tiny]left:$\blacktriangleright$}},
  aread/.style={draw=aRead, line width=1.6pt, dash pattern=on 3pt off 2pt},
  awrite/.style={draw=aWrite, line width=2.4pt},
  hw/.style={draw=nLine, thick, fill=nFill, text=nInk, rounded corners=3pt, align=center, inner sep=4pt},
  hwmain/.style={draw=nInk, thick, fill=nInk, text=white, rounded corners=3pt, align=center, inner sep=4pt},
  cell/.style={draw=nLine, thick, fill=white, text=nInk, font=\ttfamily\scriptsize},
  cellused/.style={cell, fill=nFill},
  cellfree/.style={cell, dashed, draw=nLight},
  container/.style={draw=nLine, thick, dashed, rounded corners=5pt},
  actor/.style={fill=#1, draw=#1, text=white, rounded corners=3pt, align=center, font=\scriptsize\bfseries, inner sep=4pt},
  actorlight/.style={fill=#1!15, draw=#1, thick, text=nInk, rounded corners=3pt, align=center, inner sep=4pt},
  suspended/.style={fill=#1!30, draw=#1, dashed, line width=0.9pt, text=nInk, rounded corners=3pt, align=center, inner sep=4pt},
  ready/.style={fill=white, draw=#1, line width=0.8pt, text=nInk, rounded corners=3pt, align=center, inner sep=4pt},
  idle/.style={fill=nIdle, draw=none, text=nInk, rounded corners=3pt, align=center, inner sep=4pt},
  flow/.style={-{Stealth}, very thick, DeepBlue},
  data/.style={-{Stealth}, very thick, nInk},
  holds/.style={-{Stealth}, very thick, nInk},
  waits/.style={-{Stealth}, very thick, nLine, dashed},
  annot/.style={draw=AccentOrange, line width=1.4pt, rounded corners=3pt},
}
\tikzset{actoramber/.style={actor=sAmber, text=black}}

% ---------- signs and small helpers ----------
\newcommand{\ok}{\textcolor{aExec}{\faCheck}}
\newcommand{\nok}{\textcolor{aWrite}{\faTimes}}
\newcommand{\badge}[1]{\tikz[baseline=(b.base)]\node[circle,fill=AccentOrange,text=white,
  inner sep=1.2pt,font=\bfseries\tiny] (b) {#1};}
\newcommand{\bigbadge}[1]{\tikz[baseline=(b.base)]\node[circle,fill=AccentOrange,text=white,
  inner sep=2pt,font=\bfseries\small] (b) {#1};}
% ---------- text boxes ----------
\newcommand{\defbox}[1]{%
  \begin{beamercolorbox}[rounded=true,sep=0.1cm,wd=\linewidth]{block body example}
    \footnotesize #1
  \end{beamercolorbox}\vspace{0.06cm}}
% a large definition in the middle of the slide
\newcommand{\bigdef}[1]{%
  \begin{beamercolorbox}[rounded=true,sep=0.25cm,wd=\linewidth]{block body example}
    \centering\normalsize #1
  \end{beamercolorbox}\vspace{0.06cm}}
\newcommand{\stepbox}[2]{%
  \begin{beamercolorbox}[rounded=true,sep=0.09cm,wd=\linewidth]{block body}
    \footnotesize\textbf{#1}\\ #2
  \end{beamercolorbox}\vspace{0.06cm}}
\newcommand{\keymsg}[1]{%
  \begin{beamercolorbox}[rounded=true,sep=0.1cm,wd=\textwidth]{block body example}
    \centering\footnotesize #1
  \end{beamercolorbox}}
% a grey box without a title
\newcommand{\plainbox}[1]{%
  \begin{beamercolorbox}[rounded=true,sep=0.12cm,wd=\linewidth]{block body}
    \centering\footnotesize #1
  \end{beamercolorbox}\vspace{0.06cm}}
\setbeamercolor{block body problem}{bg=AccentOrange!10,fg=black}
\newcommand{\probbox}[2]{%
  \begin{beamercolorbox}[rounded=true,sep=0.09cm,wd=\linewidth]{block body problem}
    \footnotesize\textbf{\textcolor{AccentOrange}{\faExclamationTriangle}\ #1}\\ #2
  \end{beamercolorbox}\vspace{0.06cm}}

% ---------- code: dark window, grey dots, syntax colours ----------
\definecolor{VSDarkBg}{RGB}{25,25,25}
\definecolor{VSText}{RGB}{245,245,245}
\definecolor{VSBlue}{RGB}{80,180,255}
\definecolor{VSPurple}{RGB}{210,110,210}
\definecolor{VSString}{RGB}{255,140,80}
\definecolor{VSComment}{RGB}{80,220,80}
\setbeamercolor{vscodelisting}{bg=VSDarkBg}
\lstdefinestyle{code}{basicstyle=\color{VSText}\ttfamily\scriptsize, keywordstyle=\color{VSBlue}\bfseries,
  stringstyle=\color{VSString}, commentstyle=\color{VSComment}\itshape, emphstyle=\color{VSPurple},
  numbers=none, columns=fixed, basewidth=0.52em, keepspaces=true, showstringspaces=false, xleftmargin=6pt}
\lstdefinestyle{py}{style=code, language=Python, emph={print,chr,zip,range,len,bytes,ord,hex}}
\BeforeBeginEnvironment{lstlisting}{%
  \begin{beamercolorbox}[rounded=true,sep=2mm,wd=\linewidth]{vscodelisting}%
  \begin{tikzpicture}
    \fill[nLine] (0,0) circle (3pt); \fill[nLine] (0.35,0) circle (3pt); \fill[nLine] (0.70,0) circle (3pt);
  \end{tikzpicture}%
  \vspace{-0.2cm}}
\AfterEndEnvironment{lstlisting}{\end{beamercolorbox}}
\newcommand{\hlcode}[1]{{\setlength{\fboxsep}{1pt}\colorbox{AccentOrange}{\textcolor{white}{\bfseries #1}}}}
\newcommand{\eol}{\tikz[remember picture, overlay]\coordinate (eol) at (0,0.08);}
\newcommand{\codenote}[3][0pt]{%
  \begin{tikzpicture}[remember picture, overlay, >=Stealth]
    \node[draw=AccentOrange, very thick, fill=AccentOrange!6, rounded corners=4pt, text width=0.33\paperwidth,
      align=left, inner sep=6pt, font=\footnotesize, anchor=west] (cal)
      at ([yshift=#1]eol -| {$(current page.west)!0.6!(current page.east)$}) {\textbf{\textcolor{AccentOrange}{#2}}\\[3pt]#3};
    \draw[->, very thick, AccentOrange] (eol -| cal.west) -- ([xshift=4pt]eol);
  \end{tikzpicture}}
% program output: a light terminal line
\newcommand{\terminal}[1]{%
  \begin{tikzpicture}
    \node[draw=nLine, thick, fill=nFill, rounded corners=3pt, inner sep=5pt, font=\ttfamily\small, text=nInk]
      {\textcolor{nLine}{\$}\ #1};
  \end{tikzpicture}}

% ======================================================================
% Course macros
% ======================================================================
% the starter files on GitHub: the word "GitHub" is the link.
%   \ghlink{Task4} opens the folder of Task 4, \ghlink{} opens the folder of this lecture
\newcommand{\ghbase}{https://github.com/KsaweryM/Cryptography-Course/tree/main/Lecture8}
\newcommand{\ghlink}[1]{\href{\ghbase/#1}{\textcolor{DeepBlue}{\textbf{GitHub}}}}
% the task slides: the steps are in the README on GitHub, the slide only opens the task
%   \taskgoal{one line}   \taskready{the ready modules}   \opentask{Task2}{about 6 minutes}
%   a task on paper: the question on the slide, then \onpaper{about 6 minutes}
\newcommand{\taskgoal}[1]{{\centering\large\textcolor{AccentOrange}{\faBullseye}\ \textbf{Goal:} #1\par}\vspace{0.2cm}}
\newcommand{\taskready}[1]{{\centering\normalsize\textcolor{nInk}{\faBox}\ \textbf{Ready:} #1\par}\vspace{0.2cm}}
\newcommand{\opentask}[2]{%
  \begin{beamercolorbox}[rounded=true,sep=0.35cm,wd=\linewidth]{block body example}
    \centering
    {\LARGE\faGithub\ \textbf{Open the task:}\ \href{\ghbase/#1}{\textcolor{DeepBlue}{\textbf{\texttt{#1}}}}}\\[8pt]
    \footnotesize\href{\ghbase/#1}{\textcolor{DeepBlue}{\texttt{github.com/KsaweryM/Cryptography-Course/tree/main/Lecture8/#1}}}\\[6pt]
    \small on \ghlink{#1}: the steps, the starter files and how to run them\ifx&#2&\else\qquad\faClock[regular]\ #2\fi
  \end{beamercolorbox}}
\newcommand{\onpaper}[1]{%
  \begin{beamercolorbox}[rounded=true,sep=0.3cm,wd=\linewidth]{block body example}
    \centering
    {\LARGE\faPencil*\ \textbf{On paper}}\ifx&#1&\else\qquad\small\faClock[regular]\ #1\fi
  \end{beamercolorbox}}

% a page of the simulator website, printed in full so that students can type it:
%   \simlink{elliptic-curves/\#2} prints and links ksawerym.github.io/Cryptography-simulators/elliptic-curves/#2
\newcommand{\simbase}{ksawerym.github.io/Cryptography-simulators}
\newcommand{\simlink}[1]{\href{https://\simbase/#1}{\textcolor{DeepBlue}{\texttt{\simbase/#1}}}}
% a slide that sends the room to the simulator (shown live by the lecturer): \opensim{elliptic-curves/\#1}
\newcommand{\opensim}[1]{%
  \begin{beamercolorbox}[rounded=true,sep=0.35cm,wd=\linewidth]{block body example}
    \centering
    {\LARGE\faDesktop\ \textbf{Open the simulator}}\\[8pt]
    \normalsize\simlink{#1}
  \end{beamercolorbox}}
% "For the curious" line under a frame: \curious{elliptic-curves/\#2}
\newcommand{\curious}[1]{{\footnotesize\textcolor{gray}{\faSearch\ For the curious:}\ \simlink{#1}}}

% speaker notes (Polish): printed only in the notes PDF (make notes)
\setbeamertemplate{note page}[default]
\setbeamerfont{note page}{size=\normalsize}
\newcommand{\say}[1]{\note{\selectlanguage{polish}#1}}

% skeleton macros (kept for later passes; not used in the final slides)
\newcommand{\sk}[1]{{\centering\small\textcolor{nInk}{\faPen\ \itshape #1}\par}\vspace{0.2cm}}
\newcommand{\fig}[2][3.2cm]{{\centering\fcolorbox{nLine}{nFill}{\parbox[c][#1][c]{0.8\linewidth}{\centering\small\textcolor{nInk}{\faImage[regular]\ #2}}}\par}\vspace{0.2cm}}

% ======================================================================
% Colour mapping for Lectures 6 to 8 (colour rules, section 6); the same in the task figures and the simulator
%   S1 sIndigo : Alice, and what only Alice can make: her private key (full), her signature (full, with a seal);
%                her public key is light indigo
%   S2 sPink   : Eve and Mallory, the attackers (Mallory's keys and signatures too)
%   S3 sViolet : the message, the data, the file (the input)
%   S4 sAmber  : a secret shared value: the MAC key of Alice and Bob
%   S5 sTeal   : Bob, the verifier, and Bob's key pair (public key light, private key full, as in Lecture 3)
%   S6 sBrown  : the digest (hash value) and the MAC tag: the fixed-size output that travels or is stored
%   mechanisms (SHA-256, HMAC, sign, verify, RSA, the network) are neutral dark; Carol and Dave are neutral
%   results: \ok green, \nok red
% ======================================================================
\tikzset{
  person/.style={circle, fill=#1, draw=#1, text=white, minimum size=1.1cm, font=\Large, inner sep=0pt},
  pname/.style={font=\small\bfseries, text=nInk},
  msg/.style={draw=sViolet, thick, fill=sViolet!15, text=nInk, rounded corners=3pt, align=center, inner sep=4pt},
  digest/.style={draw=sBrown, thick, fill=sBrown!15, text=nInk, rounded corners=3pt, align=center, inner sep=4pt},
  mackey/.style={draw=sAmber, thick, fill=sAmber!30, text=nInk, rounded corners=3pt, align=center, inner sep=4pt},
  apriv/.style={draw=sIndigo, thick, fill=sIndigo, text=white, rounded corners=3pt, align=center, inner sep=4pt},
  apub/.style={draw=sIndigo, thick, fill=sIndigo!15, text=nInk, rounded corners=3pt, align=center, inner sep=4pt},
  asig/.style={draw=sIndigo, thick, fill=sIndigo, text=white, rounded corners=3pt, align=center, inner sep=4pt},
  bpriv/.style={draw=sTeal, thick, fill=sTeal, text=white, rounded corners=3pt, align=center, inner sep=4pt},
  bpub/.style={draw=sTeal, thick, fill=sTeal!15, text=nInk, rounded corners=3pt, align=center, inner sep=4pt},
  mpriv/.style={draw=sPink, thick, fill=sPink, text=white, rounded corners=3pt, align=center, inner sep=4pt},
  mpub/.style={draw=sPink, thick, fill=sPink!15, text=nInk, rounded corners=3pt, align=center, inner sep=4pt},
  msig/.style={draw=sPink, thick, fill=sPink, text=white, rounded corners=3pt, align=center, inner sep=4pt},
  capt/.style={font=\scriptsize, text=gray, align=center},
  net/.style={-{Stealth}, very thick, nInk, decorate, decoration={snake, amplitude=1.4pt, segment length=8pt, post length=4pt}},
}
% people: \alice{pos}, \bob{pos}, \mallory{pos}, \eve{pos}, \carol{pos}, \dave{pos}
% (each draws a node A, B, M, E, C or D and its name below)
\newcommand{\alice}[1]{\node[person=sIndigo] (A) at (#1) {\faUser}; \node[pname, below=1pt of A] (Albl) {Alice};}
\newcommand{\bob}[1]{\node[person=sTeal] (B) at (#1) {\faUser}; \node[pname, below=1pt of B] (Blbl) {Bob};}
\newcommand{\mallory}[1]{\node[person=sPink] (M) at (#1) {\faUserSecret}; \node[pname, below=1pt of M] (Mlbl) {Mallory};}
\newcommand{\eve}[1]{\node[person=sPink] (E) at (#1) {\faUserSecret}; \node[pname, below=1pt of E] (Elbl) {Eve};}
\newcommand{\carol}[1]{\node[person=nInk] (C) at (#1) {\faUserTie}; \node[pname, below=1pt of C] (Clbl) {Carol};}
\newcommand{\dave}[1]{\node[person=nInk] (D) at (#1) {\faUser}; \node[pname, below=1pt of D] (Dlbl) {Dave};}
% a mechanism box with cogs: \mech{name}{pos}{text}{width}
\newcommand{\mech}[4]{\node[hwmain, minimum width=#4, minimum height=1.0cm, font=\small\bfseries] (#1) at (#2) {\faCogs\ #3};}
% words in the colour of their element (only when the element is on the same slide)
\newcommand{\msgw}[1]{\textcolor{sViolet}{\textbf{#1}}}
\newcommand{\aliw}[1]{\textcolor{sIndigo}{\textbf{#1}}}
\newcommand{\bobw}[1]{\textcolor{sTealText}{\textbf{#1}}}
\newcommand{\malw}[1]{\textcolor{sPinkText}{\textbf{#1}}}
\newcommand{\digw}[1]{\textcolor{sBrown}{\textbf{#1}}}
\newcommand{\keyw}[1]{\textcolor{sAmberText}{\textbf{#1}}}
\newcommand{\hl}[1]{\textcolor{AccentOrange}{\textbf{#1}}}
% a step of a calculation: an orange arrow when it is the step of this slide, nothing before it otherwise
\newcommand{\now}{\makebox[0.45cm][l]{\textcolor{AccentOrange}{\faCaretRight}}}
\newcommand{\notnow}{\makebox[0.45cm][l]{}}

% ---------- the toy curve y^2 = x^3 + 2x + 2 mod 17 ----------
% CHECKED: python3, all (x, y) with 0 <= x, y < 17 and y^2 = x^3 + 2x + 2 mod 17: these 18 points
\newcommand{\curvepts}{0/6,0/11,3/1,3/16,5/1,5/16,6/3,6/14,7/6,7/11,9/1,9/16,10/6,10/11,13/7,13/10,16/4,16/13}
% the grid of dots: \modgrid[cell size in cm]{more drawing, in grid coordinates}
\newcommand{\modgrid}[2][0.3]{%
  \begin{tikzpicture}[x=#1cm, y=#1cm, >=Stealth]
    \foreach \x in {0,...,16} \foreach \y in {0,...,16} \fill[nLight] (\x,\y) circle (0.9pt);
    \draw[nLine, thick] (-0.6,-0.6) -- (16.6,-0.6);
    \draw[nLine, thick] (-0.6,-0.6) -- (-0.6,16.6);
    \foreach \v in {0,...,16} {
      \node[font=\tiny, text=gray, anchor=north] at (\v,-0.65) {\v};
      \node[font=\tiny, text=gray, anchor=east] at (-0.65,\v) {\v};
    }
    \node[font=\scriptsize, text=nInk, anchor=west] at (16.7,-0.6) {$x$};
    \node[font=\scriptsize, text=nInk, anchor=south] at (-0.6,16.7) {$y$};
    \foreach \x/\y in \curvepts \fill[nInk] (\x,\y) circle (2.6pt);
    #2
  \end{tikzpicture}}
% mark a point of the grid: \gmark[colour]{x,y}{label}{where the label goes}
\newcommand{\gmark}[4][sViolet]{%
  \node[circle, draw=#1, line width=1.2pt, fill=#1, minimum size=8pt, inner sep=0pt] at (#2) {};
  \node[font=\scriptsize\bfseries, text=#1, #4, inner sep=1pt, outer sep=3pt, fill=white, fill opacity=0.85,
    text opacity=1, rounded corners=1pt] at (#2) {#3};}
% an empty place for a point that is not known yet
\newcommand{\gempty}[2]{%
  \node[circle, draw=AccentOrange, dashed, line width=1.2pt, minimum size=11pt, inner sep=0pt] at (#1) {};
  \node[font=\scriptsize\bfseries, text=AccentOrange, #2] at (#1) {?};}


\title{Cryptography}
\subtitle{Digital Signatures}
\author{Ksawery Możdżyński}
\date{}

\begin{document}

% ======================================================================
% Opening
% ======================================================================

\begin{frame}
  \titlepage
  \say{Dzisiejsze pytanie: jak udowodnić, kto napisał wiadomość? Na wykładach 6 i 7 skrót mówił, czy dane się zmieniły. Dziś zobaczymy, że to za mało, i poznamy podpis cyfrowy. Cały wykład trwa około 4 godzin razem z 9 zadaniami (jedno na kartce). Zadania 2 do 9 potrzebują biblioteki cryptography (pip install cryptography), tej samej co w wykładach 9 do 12.}
\end{frame}

\begin{frame}{Agenda For Today}
  \vfill
  \begin{center}
    \begin{tikzpicture}
      \foreach \ic/\t [count=\k from 0] in {
          {\faUserSecret}/{See why a hash or a shared key cannot prove who wrote a message, and how a signature does.},
          {\faCogs}/{Sign with RSA by hand, forge such a signature without a hash, and fix it by hashing first.},
          {\faBalanceScale}/{Compare signatures with MACs, break ECDSA when its random number repeats, and see where signatures are used.}} {
        \node[circle, fill=DeepBlue, text=white, minimum size=1.3cm, font=\LARGE] at (0,{-1.9*\k}) {\ic};
        \node[anchor=west, text width=12cm, font=\large, align=left, execute at begin node={\hyphenpenalty=10000 }] at (1.0,{-1.9*\k}) {\t};
      }
    \end{tikzpicture}
  \end{center}
  \vfill
  \say{Trzy części: najpierw problem i pomysł podpisu (z gotowym podpisem Ed25519), potem zajrzymy do środka podpisu RSA i jego pułapek, na końcu porównanie z MAC, atak na ECDSA i zastosowania. Nazwy sekcji są na pasku u góry.}
\end{frame}

% ======================================================================
\section[Who wrote this?]{Who wrote this?}
% ======================================================================

\begin{frame}{The update of Alice's app}
  \vfill
  \begin{center}
    \begin{tikzpicture}[>=Stealth]
      \alice{0,0}
      % the server with the download page
      \node[hw, minimum width=5.6cm, minimum height=3.3cm] (srv) at (4.6,0) {};
      \node[font=\small\bfseries, text=nInk, anchor=north west] at ([xshift=4pt,yshift=-3pt]srv.north west) {\faServer\ notes.example/download};
      \node[msg, minimum width=4.8cm, font=\small] (file) at (4.6,0.35) {\faFile*[regular]\ \texttt{update.sh}\ \ Notes 2.1};
      \node[digest, minimum width=4.8cm, font=\scriptsize\ttfamily] (h) at (4.6,-0.85) {SHA-256: ae820a2ebc8d26f1\ldots 32b27fb4};
      \draw[data] (A.east) to node[above, font=\scriptsize] {publishes} (srv.west);
      \bob{10.4,0}
      \draw[net] (srv.east) to node[above, font=\scriptsize] {downloads} (B.west);
    \end{tikzpicture}
  \end{center}
  \vfill
  \begin{center}\small
    Alice writes the Notes app. She puts each update and its \digw{SHA-256} on her download page.\\[4pt]
    Bob installs the \msgw{update}. It runs on his computer with all his rights.
  \end{center}
  \vfill
  \keymsg{Before he runs it, Bob wants to know: did Alice write this file?}
  % CHECKED: python3 hashlib.sha256 of Task1/alice_update.sh = ae820a2ebc8d26f19855ff947c6054973ac06b4c734dc966a9b1a98b32b27fb4
  \say{Konkretna sytuacja na cały pierwszy blok: aktualizacja programu. To bardzo częsty przypadek: aktualizacja ma pełne prawa na komputerze, więc fałszywa aktualizacja to najlepszy prezent dla atakującego. Skrót na slajdzie to prawdziwy SHA-256 pliku alice\_update.sh z zadania 1. Zapytać salę: co Bob może zrobić, żeby sprawdzić plik? (Oczekiwana odpowiedź: porównać skrót.)}
\end{frame}

\samedot
\begin{frame}{Recall: a hash function}
  % CHECKED: python3 hashlib.sha256(b"Pay Bob 10 pounds.") and (b"Pay Bob 90 pounds.")
  \vfill
  \begin{center}
    \begin{tikzpicture}[>=Stealth]
      \foreach \y/\in/\out in {0.8/{Pay Bob 10 pounds.}/{ded9d3862c934a76ee6fb4f57ebd18525f31f2b11a64e05c0e16a255875f149e},
          -0.8/{Pay Bob 90 pounds.}/{3afa0d2322d913e8a19c35cf5e6c6f25a79dfecc47d815812fa6fc4ab1703922}} {
        \node[msg, font=\small\ttfamily, minimum width=3.2cm] (i) at (0,\y) {\in};
        \node[hwmain, minimum width=2.0cm, minimum height=0.9cm, font=\small\bfseries] (h) at (3.3,\y) {\faCogs\ SHA-256};
        \node[digest, font=\tiny\ttfamily, anchor=west, minimum height=0.6cm] (o) at (4.8,\y) {\out};
        \draw[data] (i) -- (h);
        \draw[data] (h) -- (o);
      }
      \node[font=\scriptsize, text=AccentOrange, anchor=north] at (0,0.3) {one digit changed};
      \node[font=\scriptsize, text=gray, anchor=north] at (9.9,-1.2) {always 32 bytes, completely different};
    \end{tikzpicture}
  \end{center}
  \vfill
  \begin{columns}[t,onlytextwidth]
    \column{0.32\textwidth}\stepbox{Fixed size}{Any input gives a 32-byte \digw{digest}.}
    \column{0.32\textwidth}\stepbox{One way}{From the digest you cannot get the input back.}
    \column{0.32\textwidth}\stepbox{No collisions}{Nobody can find two inputs with the same digest.}
  \end{columns}
  \say{Przypomnienie z wykładu 6. Digest to wynik funkcji skrótu (po polsku: skrót). Trzy własności: stały rozmiar, jednokierunkowość (nie da się znaleźć wejścia do danego skrótu: preimage), brak kolizji (wykład 7). Dziś skrót będzie potrzebny dwa razy: przy sprawdzaniu pobranego pliku i w środku podpisu. Skróty policzone w Pythonie, hashlib.sha256.}
\end{frame}

\samedot
\begin{frame}{Bob checks the download}
  \vfill
  \begin{center}
    \begin{tikzpicture}[>=Stealth]
      \bob{0,0}
      \node[msg, font=\small, minimum width=3.0cm] (f) at (3.2,0.9) {\faFile*[regular]\ \texttt{update.sh}};
      \mech{h}{6.4,0.9}{SHA-256}{2.4cm}
      \node[digest, font=\scriptsize\ttfamily] (d1) at (9.9,0.9) {ae820a2e\ldots 32b27fb4};
      \node[font=\scriptsize, text=gray, anchor=south] at (d1.north) {computed by Bob};
      \node[digest, font=\scriptsize\ttfamily] (d2) at (9.9,-0.9) {ae820a2e\ldots 32b27fb4};
      \node[font=\scriptsize, text=gray, anchor=north] at (d2.south) {read on the download page};
      \draw[data] (B.east) -- (f.west);
      \draw[data] (f) -- (h);
      \draw[data] (h) -- (d1);
      \node[font=\Large] at (12.1,0) {= \ok};
    \end{tikzpicture}
  \end{center}
  \vfill
  \begin{columns}[t,onlytextwidth]
    \column{0.48\textwidth}\stepbox{\ok\ This catches}{a broken download; a file changed on a mirror server that does not host the page.}
    \column{0.48\textwidth}\stepbox{The idea}{one changed byte gives a completely different digest.}
  \end{columns}
  \keymsg{Equal digests: the file is the one the page describes.}
  \say{Bob liczy skrót sam i porównuje go ze skrótem ze strony. Serwer lustrzany (mirror) to kopia plików na innym serwerze, np. bliżej użytkownika; strona ze skrótem jest u Alice, plik na mirrorze. Tak naprawdę robi się to dla obrazów systemów (np. pliki SHA256SUMS przy obrazach Linuksa). Podkreślić słowa ,,the page describes'': skrót porównuje plik ze stroną, a nie z Alice.}
\end{frame}

\samedot
\begin{frame}{Mallory replaces both}
  \vfill
  \begin{center}
    \begin{tikzpicture}[>=Stealth]
      \mallory{0,0}
      \node[hw, minimum width=5.6cm, minimum height=3.3cm] (srv) at (4.6,0) {};
      \node[font=\small\bfseries, text=nInk, anchor=north west] at ([xshift=4pt,yshift=-3pt]srv.north west) {\faServer\ notes.example/download};
      \node[msg, minimum width=4.8cm, font=\small, draw=AccentOrange, line width=1.4pt] (file) at (4.6,0.35) {\faFile*[regular]\ \texttt{update.sh} + one line of Mallory};
      \node[digest, minimum width=4.8cm, font=\scriptsize\ttfamily, draw=AccentOrange, line width=1.4pt] (h) at (4.6,-0.85) {SHA-256: b4d421a73b76fa85\ldots 7aad982d};
      \draw[data] (M.east) to node[above, font=\scriptsize] {changes both} (srv.west);
      \bob{10.4,0}
      \draw[net] (srv.east) -- (B.west);
      \node[font=\small, text=nInk, anchor=north] at (10.4,-1.3) {digest check: \ok};
    \end{tikzpicture}
  \end{center}
  \vfill
  \begin{center}\small
    Mallory breaks into the server (or sits between Bob and it). She adds one line to the \msgw{update}\\
    and puts the \digw{SHA-256} of her file on the page. Bob's check passes.
  \end{center}
  \keymsg{Anybody can compute a SHA-256, Mallory too: a digest says nothing about who made the file.}
  % CHECKED: python3 hashlib.sha256 of Task1/mallory_update.sh = b4d421a73b76fa85cff1735acc3396eff57390b3f512e36115c827a87aad982d
  \say{To jest dokładnie przypadek z wykładu 6 (,,what a hash alone does not give''). Skrót nie ma żadnego sekretu: Mallory liczy go tak samo jak Alice. Prawdziwe przykłady włamań na serwer pobierania z podmianą pliku i sumy kontrolnej się zdarzały; szczegółów nie podajemy, chodzi o zasadę. Linia Mallory w zadaniu 1: curl ... | sh, czyli pobierz i uruchom program Mallory.}
\end{frame}

\begin{frame}{What Bob needs}
  \vfill
  \bigdef{A value that \textbf{only Alice can make} for this file, and that \textbf{Bob can check}.}
  \vfill
  \begin{columns}[t,onlytextwidth]
    \column{0.48\textwidth}
      \centering
      \begin{tikzpicture}
        \node[mackey, font=\Large, minimum width=1.6cm, minimum height=1.1cm] {\faKey};
      \end{tikzpicture}\\[4pt]
      \textbf{Way 1: a shared secret key}\\[2pt]
      \small Alice and Bob agree on one key; Mallory does not know it.
    \column{0.48\textwidth}
      \centering
      \begin{tikzpicture}
        \node[apriv, font=\Large, minimum width=1.2cm, minimum height=1.1cm] at (0,0) {\faKey};
        \node[apub, font=\Large, minimum width=1.2cm, minimum height=1.1cm] at (1.5,0) {\faKey};
      \end{tikzpicture}\\[4pt]
      \textbf{Way 2: a key pair}\\[2pt]
      \small Alice has a private key; everybody has her public key.
  \end{columns}
  \vfill
  \keymsg{The value must depend on a secret that Mallory does not have.}
  \say{Wniosek z poprzedniego slajdu: brakuje sekretu. Dwie drogi: wspólny tajny klucz (droga 1, zaraz) albo para kluczy (droga 2, po zadaniu 1). Kolory: bursztynowy klucz to wspólny sekret Alice i Boba, indygo to klucze Alice (pełny: prywatny, jasny: publiczny).}
\end{frame}

\begin{frame}{Recall: the tag of ChaCha20-Poly1305}
  \vfill
  \begin{center}
    \begin{tikzpicture}[>=Stealth]
      \node[msg, font=\small, minimum width=2.6cm] (m) at (0,0.6) {ciphertext};
      \node[mackey, font=\small, minimum width=2.6cm] (k) at (0,-0.6) {\faKey\ shared key};
      \mech{p}{3.8,0}{Poly1305}{2.6cm}
      \node[digest, font=\small, minimum width=2.2cm] (t) at (7.2,0) {tag, 16 bytes};
      \draw[data] (m.east) -- (p.west |- m.east);
      \draw[data] (k.east) -- (p.west |- k.east);
      \draw[data] (p) -- (t);
      \node[font=\small, text=nInk, anchor=west, align=left] at (8.6,0) {the receiver computes\\it again and compares};
    \end{tikzpicture}
  \end{center}
  \vfill
  \begin{center}\small
    In Lecture 2, the tag stopped Mallory from flipping bits of a ciphertext:\\
    without the key she cannot compute the right tag for her changed ciphertext.
  \end{center}
  \keymsg{A tag made with a secret key: only key holders can make it.}
  \say{Przypomnienie z wykładu 2: ChaCha20-Poly1305 dokleja do szyfrogramu 16-bajtowy tag. Odbiorca liczy tag sam i porównuje; jeśli różny, odrzuca wiadomość. Mallory mogła odwracać bity szyfrogramu (bit flipping z wykładu 1), ale nie umie policzyć nowego tagu. To jest właśnie droga 1. Teraz ta sama idea bez szyfrowania, dla dowolnej wiadomości: MAC.}
\end{frame}

\samedot
\begin{frame}{A MAC: a tag made with a shared key}
  \vfill
  \begin{center}
    \begin{tikzpicture}[>=Stealth]
      \node[msg, font=\small, minimum width=2.6cm] (m) at (0,0.6) {\faFile*[regular]\ message};
      \node[mackey, font=\small, minimum width=2.6cm] (k) at (0,-0.6) {\faKey\ shared key};
      \mech{p}{3.8,0}{HMAC-SHA256}{3.0cm}
      \node[digest, font=\small, minimum width=2.2cm] (t) at (7.2,0) {tag, 32 bytes};
      \draw[data] (m.east) -- (p.west |- m.east);
      \draw[data] (k.east) -- (p.west |- k.east);
      \draw[data] (p) -- (t);
    \end{tikzpicture}
  \end{center}
  \vfill
  \defbox{\textbf{MAC} (message authentication code): a tag computed from the message and a secret key that sender and receiver share.
    The receiver computes the tag again and compares. \textbf{HMAC-SHA256}: a MAC built from SHA-256, which mixes the key in twice.}
  \probbox{Why not simply SHA-256(key $\|$ message)?}{SHA-256 works block after block (Lecture 6). From that digest alone, Mallory can continue the hash and get a valid tag for message $\|$ her extra bytes, without the key (\textbf{length extension}). HMAC is built so that this does not work.}
  \say{Jeden slajd o MAC. HMAC (RFC 2104, 1997) liczy SHA-256 dwa razy: najpierw skrót z klucza (zmieszanego ze stałą) i wiadomości, potem skrót z klucza (zmieszanego z inną stałą) i tego pierwszego skrótu. Nie trzeba tego pamiętać; ważne, że w Pythonie jest to jedna funkcja: hmac.new(key, message, "sha256"). Znak $\|$ to sklejenie bajtów. Length extension: wynik SHA-256 to cały stan po ostatnim bloku (Merkle-Damgård z wykładu 6), więc Mallory może z niego ruszyć dalej i dopisać bloki; nie zna klucza, ale nie musi. Gdyby ktoś pytał: SHA-3 nie ma tego problemu. MAC jest wszędzie, gdzie obie strony mają wspólny klucz: po uzgodnieniu klucza w protokole przeglądarki (wykład 4), w ciasteczkach serwera, w tokenach API.}
\end{frame}

\begin{frame}{Task 1: Mallory and the tag}
  \vfill
  \taskgoal{as Bob, check Mallory's update with the digest and with the tag; then make a tag yourself.}
  \taskready{\texttt{hashlib} and \texttt{hmac} from the standard library; the two update files.}
  \vfill
  \opentask{Task1}{about 10 minutes}
  \vfill
  \say{Około 10 minut. Polecenia w README na GitHubie, tylko biblioteka standardowa. TODO 1: skrót pliku Mallory zgadza się ze stroną. TODO 2: tag (HMAC) pliku Alice się zgadza, pliku Mallory nie. TODO 3: Bob sam robi tag dla notatki ,,Alice owes Bob 500 pounds''. Przed TODO 3 zapytać: czy Carol, która dostanie klucz, odróżni notatkę od prawdziwej wiadomości Alice? Jeśli ktoś utknie: hmac.new(SHARED\_KEY, data, "sha256").digest() i hmac.compare\_digest(a, b).}
\end{frame}

\samedot
\begin{frame}{Task 1: what we saw}
  % CHECKED: python3 Task1/solution/who_made_it.py, 2026-11-17: True, True, False, True
  \vfill
  \begin{center}
    \renewcommand{\arraystretch}{1.5}
    \begin{tabular}{@{}l l c@{}}
      \toprule
      \textbf{check} & \textbf{of} & \textbf{result} \\
      \midrule
      digest from the page & Mallory's update & \ok\ \texttt{True} \\
      tag with the shared key & Alice's update & \ok\ \texttt{True} \\
      tag with the shared key & Mallory's update & \nok\ \texttt{False} \\
      tag with the shared key & Bob's note ``Alice owes Bob 500 pounds'' & \ok\ \texttt{True} \\
      \bottomrule
    \end{tabular}
  \end{center}
  \vfill
  \begin{center}\small The tag stopped Mallory. But the note that Bob wrote himself passed Carol's check.\end{center}
  \keymsg{A tag proves ``made by someone with the key'', and Bob has the key too.}
  \say{Wynik sprawdzony rozwiązaniem. Wiersz 1: skrót przepuszcza podmianę. Wiersz 3: Mallory nie ma klucza, więc jej plik nie ma poprawnego tagu. Wiersz 4 to najważniejsze: Bob zrobił poprawny tag dla notatki, której Alice nigdy nie napisała. Tag świadczy tylko o tym, że zrobił go ktoś z kluczem. Zapytać salę: kto ma ten klucz? (Alice i Bob.)}
\end{frame}

\begin{frame}{Who could have made the tag?}
  \vfill
  \begin{center}
    \begin{tikzpicture}[>=Stealth]
      \alice{0,0}
      \bob{4.4,0}
      \node[mackey, font=\small] at (0,-1.6) {\faKey\ the key};
      \node[mackey, font=\small] at (4.4,-1.6) {\faKey\ the key};
      \carol{10.4,0}
      \node[msg, font=\small, text width=3.4cm, align=center] (n) at (7.4,1.2) {``Alice owes Bob 500 pounds.''};
      \node[digest, font=\small] (t) at (7.4,0.1) {tag \ok};
      \node[font=\Huge\bfseries, text=AccentOrange] at (10.4,1.5) {?};
      \draw[data] (n.east) -- (C.north west);
      \draw[data] (t.east) -- (C.west);
    \end{tikzpicture}
  \end{center}
  \vfill
  \begin{center}\small
    Bob takes the note to Carol, a judge. The tag is valid. But Alice and Bob have the same key,\\
    so Carol cannot tell whether Alice or Bob made it.
  \end{center}
  \keymsg{With a shared key, the tag convinces Bob, but nobody else.}
  \say{To jest pierwsza słabość MAC: nie da się nim przekonać osoby trzeciej. Bob wie, że to nie on zrobił tag (jeśli mówi prawdę), więc sam jest przekonany. Ale Carol, sędzia, nie ma jak odróżnić. Zaznaczyć: to nie zawsze wada. Wrócimy do tego w sekcji 3 (kiedy MAC jest lepszy).}
\end{frame}

\samedot
\begin{frame}{Every checker can also forge}
  \vfill
  \begin{center}
    \begin{tikzpicture}[>=Stealth]
      \alice{0,0}
      \node[mackey, font=\small] at (0,-1.6) {\faKey};
      \foreach \i in {0,...,9} {
        \node[circle, fill=nLine, text=white, minimum size=0.6cm, font=\small, inner sep=0pt] (u\i) at ({4.0+0.85*\i},0.9) {\faUser};
        \node[mackey, font=\tiny, inner sep=1.5pt] at ({4.0+0.85*\i},0.2) {\faKey};
      }
      \node[font=\small, text=nInk] at (12.6,0.9) {\ldots};
      \draw[data] (A.east) -- (3.6,0.9);
      \node[font=\small, text=nInk] at (8.1,-0.6) {a million users of the Notes app};
      \node[font=\small, text=aWrite] at (8.1,-1.2) {each one could make a ``valid'' update for all the others};
    \end{tikzpicture}
  \end{center}
  \vfill
  \begin{center}\small
    To check a tag, you need the key. Who has the key can also make tags.\\
    For one message to many receivers, a shared key does not work.
  \end{center}
  \keymsg{A shared key gives the power to check and the power to make, together.}
  \say{Druga słabość: aktualizację sprawdza milion komputerów. Każdy musiałby mieć klucz, więc każdy mógłby zrobić fałszywą aktualizację dla innych, a klucz z miliona komputerów na pewno wycieknie. Potrzebujemy rozdzielić dwie rzeczy: robienie (tylko Alice) i sprawdzanie (wszyscy). Znamy już taki rozdział: para kluczy z wykładu 3.}
\end{frame}

\begin{frame}{Recall: a key pair}
  \vfill
  \begin{center}
    \begin{tikzpicture}[>=Stealth]
      \bob{0,0}
      \node[bpub, minimum width=4.6cm, minimum height=1.25cm, anchor=west, font=\small] (pk) at (2.0,0.85)
        {\faKey\ \textbf{public key}\\ everybody may know it};
      \node[bpriv, minimum width=4.6cm, minimum height=1.25cm, anchor=west, font=\small] (sk) at (2.0,-0.85)
        {\faKey\ \textbf{private key}\\ only Bob knows it};
      \draw[data] (B.east) -- (pk.west);
      \draw[data] (B.east) -- (sk.west);
      \node[font=\small, text=nInk, anchor=west] at (7.1,0.85) {anybody \textbf{encrypts} a message for Bob};
      \node[font=\small, text=nInk, anchor=west] at (7.1,-0.85) {only Bob \textbf{decrypts} it};
    \end{tikzpicture}
  \end{center}
  \vfill
  \begin{center}\small RSA (Lecture 3), ElGamal (Lectures 4 and 5): the public key cannot give the private key.\end{center}
  \keymsg{Encryption: the public key locks, the private key opens.}
  \say{Przypomnienie z wykładu 3. Klucze Boba w kolorze Boba: jasny publiczny, pełny prywatny. Z klucza publicznego nie da się policzyć prywatnego (w RSA trzeba by rozłożyć n). Zapytać: co by było, gdybyśmy zamienili role kluczy?}
\end{frame}

\samedot
\begin{frame}{The idea: the private key signs}
  \vfill
  \begin{center}
    \begin{tikzpicture}[>=Stealth]
      \alice{0,0}
      \node[msg, font=\small] (m) at (2.6,0.7) {\faFile*[regular]\ message};
      \node[apriv, font=\small] (sk) at (2.6,-0.7) {\faKey\ private key};
      \mech{s}{5.3,0}{sign}{1.8cm}
      \node[asig, font=\small, minimum height=0.9cm] (sig) at (7.8,0) {\faStamp\ signature};
      \draw[data] (m.east) -- (s.west |- m.east);
      \draw[data] (sk.east) -- (s.west |- sk.east);
      \draw[data] (s) -- (sig);
      \mech{v}{10.4,0}{verify}{1.8cm}
      \node[apub, font=\small] (pk) at (10.4,1.3) {\faKey\ Alice's public key};
      \draw[data] (sig) -- (v);
      \draw[data] (pk) -- (v);
      \node[font=\Large] (r) at (12.4,0) {\ok\,/\,\nok};
      \draw[data] (v) -- (r);
      \bob{10.4,-1.7}
    \end{tikzpicture}
  \end{center}
  \vfill
  \begin{center}\small
    Turn the pair around: the \aliw{private key} makes the value, the public key checks it.\\
    Only Alice can make it. Bob, Carol and a million users can check it, and none of them can make it.
  \end{center}
  \keymsg{Making and checking are now split between two keys.}
  \say{To jest cały pomysł podpisu cyfrowego (Diffie i Hellman opisali go w 1976 roku, pierwszy działający schemat to RSA w 1977). Kolory: podpis Alice to pełne indygo z pieczątką, bo tylko Alice może go zrobić; klucz publiczny Alice jasne indygo. verify dostaje też wiadomość (na rysunku pominięta, jest na następnych slajdach).}
\end{frame}

\samedot
\begin{frame}{A digital signature}
  \vfill
  \bigdef{\textbf{Digital signature}: a value made from a message with a \textbf{private key}.\\
    Anyone with the matching \textbf{public key} can check that it fits this message.}
  \vfill
  \begin{columns}[c,onlytextwidth]
    \column{0.30\textwidth}
      \centering
      \begin{tikzpicture}
        \node[bpub, font=\Large, minimum width=1.8cm, minimum height=1.1cm] at (0,0) {$B$};
        \node[circle, draw=sTeal, very thick, fill=sTeal!15, text=sTeal, minimum size=1.0cm, inner sep=0pt, font=\tiny\bfseries, align=center] at (1.6,0) {\faStamp\\Bob};
      \end{tikzpicture}
    \column{0.66\textwidth}
      \small Lecture 4, ``The fix: authenticate the exchange'': Bob added to his Diffie-Hellman value $B$ a \textbf{seal}
      that only his long-term private key can make, and Alice checked it with his public key.\\[4pt]
      \textbf{That seal is a digital signature.} Today: how it works.
  \end{columns}
  \vfill
  \keymsg{A signature: made with the private key, checked with the public key.}
  \say{Definicja. Nawiązanie do wykładu 4: tam obiecaliśmy, że ,,pieczęć'' Boba wyjaśnimy na wykładzie 8. Mallory w ataku man in the middle mogła podmienić B, ale nie umiała zapieczętować swojej wartości jako Bob. Pieczęć Boba jest w kolorze Boba, bo to jego klucz; dziś podpisuje głównie Alice. Polska nazwa: podpis cyfrowy, podpis elektroniczny to szersze pojęcie prawne.}
\end{frame}

\begin{frame}{sign and verify}
  \vfill
  \begin{center}
    \begin{tikzpicture}[>=Stealth]
      % sign
      \node[msg, font=\small] (d1) at (0,1.0) {\faFile*[regular]\ message};
      \node[apriv, font=\small] (sk) at (0,0.1) {\faKey\ Alice's private key};
      \node[hwmain, minimum width=2.0cm, minimum height=1.6cm] (sg) at (3.4,0.55) {\Large\faCogs\\ \small\bfseries sign};
      \node[asig, font=\small] (s1) at (6.2,0.55) {\faStamp\ signature};
      \draw[data] (d1.east) -- (sg.west |- d1.east);
      \draw[data] (sk.east) -- (sg.west |- sk.east);
      \draw[data] (sg) -- (s1);
      % verify
      \node[msg, font=\small] (d2) at (0,-1.6) {\faFile*[regular]\ message};
      \node[apub, font=\small] (pk) at (0,-2.4) {\faKey\ Alice's public key};
      \node[asig, font=\small] (s2) at (0,-3.2) {\faStamp\ signature};
      \node[hwmain, minimum width=2.0cm, minimum height=2.0cm] (vf) at (3.4,-2.4) {\Large\faCogs\\ \small\bfseries verify};
      \draw[data] (d2.east) -- (vf.west |- d2.east);
      \draw[data] (pk.east) -- (vf.west |- pk.east);
      \draw[data] (s2.east) -- (vf.west |- s2.east);
      \node[font=\small, anchor=west] (r) at (5.2,-2.4) {\ok\ \texttt{True}\; or \;\nok\ \texttt{False}};
      \draw[data] (vf) -- (r);
      \node[font=\small, text=nInk, anchor=west, align=left] at (8.8,0.55) {Alice, once};
      \node[font=\small, text=nInk, anchor=west, align=left] at (8.8,-1.2) {anybody, any time:\\Bob, Carol, a million users};
    \end{tikzpicture}
  \end{center}
  \vfill
  \keymsg{sign needs the private key; verify needs only the public key.}
  \say{Czarna skrzynka, bez środka. Te same nazwy (sign, verify) będą w zadaniu 2 i w wykładach 9 do 12. Uwaga: verify dostaje wiadomość, a nie tylko podpis; podpis bez wiadomości nic nie znaczy. Podpis jest osobnym kawałkiem bajtów, który podróżuje obok wiadomości (jak tag MAC).}
\end{frame}

\samedot
\begin{frame}{What True means}
  \vfill
  \begin{center}
    \begin{tikzpicture}
      \node[hw, text width=10.6cm, align=left, font=\normalsize, inner sep=10pt] {
        \textbf{\texttt{verify} says \texttt{True}:}\\[8pt]
        \bigbadge{1}\ the message was \textbf{not changed} since it was signed;\\[6pt]
        \bigbadge{2}\ it was signed with the \textbf{private key that belongs to this public key}.};
    \end{tikzpicture}
  \end{center}
  \vfill
  \begin{center}\small
    \textbf{False} says only that something does not fit: a changed message, a changed signature, or another key.\\[4pt]
    Point 2 is about a \textbf{key}, not a person. That matters at the end of today.
  \end{center}
  \keymsg{A valid signature binds this message to this public key.}
  \say{Dwa punkty, te same co w przypomnieniu na wykładzie 9. Punkt 2 mówi o kluczu: podpis udowadnia, że podpisał ktoś, kto ma klucz prywatny pasujący do tego klucza publicznego. Czy to Alice? Tylko jeśli klucz publiczny naprawdę jest Alice. Do tego wrócimy w zadaniu 9.}
\end{frame}

\begin{frame}{Ed25519}
  \vfill
  \begin{center}
    \begin{tikzpicture}
      \foreach \x/\sty/\name/\sz/\w in {0/apriv/private key/32 bytes/2.4, 4.4/apub/public key/32 bytes/2.4, 8.8/asig/signature/64 bytes/4.8} {
        \node[\sty, minimum width=\w cm, minimum height=1.0cm, font=\small] at (\x+\w/2,0) {\name};
        \node[font=\small\bfseries, text=nInk] at (\x+\w/2,-0.85) {\sz};
      }
    \end{tikzpicture}
  \end{center}
  \vfill
  \begin{itemize}
    \item A signature scheme on \textbf{Curve25519}, the curve of X25519 (Lecture 5). Bernstein and others, 2011; RFC 8032, 2017.
    \item Small and fast: the whole signature is 64 bytes.
    \item \textbf{Deterministic}: the same message and the same key always give the same signature.
  \end{itemize}
  \vfill
  \keymsg{Ed25519: 32-byte keys, 64-byte signatures, no randomness needed.}
  % CHECKED: RFC 8032 section 5.1.5 and 5.1.6: 32-byte private and public keys, 64-byte signature; Bernstein, Duif, Lange, Schwabe, Yang, "High-speed high-security signatures", CHES 2011
  \say{Ed25519 to dziś jeden z najczęściej polecanych podpisów: SSH (klucze ssh-ed25519), podpisane commity w Git, Signal, wiele bibliotek. Na razie traktujemy go jako czarną skrzynkę; w sekcji 3 powiemy, czemu deterministyczność jest ważna. Rozmiary sprawdzone w RFC 8032 i w zadaniu 2.}
\end{frame}

\samedot
\begin{frame}[fragile]{Ed25519 in Python}
\begin{lstlisting}[style=py]
from cryptography.hazmat.primitives.asymmetric.ed25519 import Ed25519PrivateKey
from cryptography.exceptions import InvalidSignature

private_key = Ed25519PrivateKey.generate()      # Alice, once
public_key = private_key.public_key()            # Alice gives this to everybody

signature = private_key.sign(b"Meet at 10")      # 64 bytes

try:
    public_key.verify(signature, b"Meet at 10")  # Bob
    print("valid")
except InvalidSignature:
    print("NOT valid")
\end{lstlisting}
  \vfill
  \begin{center}\small
    \texttt{verify} returns nothing if the signature fits and \textbf{raises an exception} if not:\\
    like a C function that jumps to the error label, so a wrong signature can never be ignored by mistake.
  \end{center}
  \say{Biblioteka cryptography, ta sama co w wykładach 9 do 12. W zadaniu 2 nie trzeba pisać try/except: gotowy moduł signing.py zamienia to na True albo False (tak samo jak w wykładzie 9). Dlaczego wyjątek: programista, który zapomni sprawdzić wartość zwracaną, zostałby z niesprawdzonym podpisem; wyjątku nie da się przeoczyć, program się zatrzyma.}
\end{frame}

\begin{frame}{Task 2: sign and verify}
  \vfill
  \taskgoal{as Alice, sign a message with Ed25519; as Bob, check it, a changed copy and Mallory's signature.}
  \taskready{\texttt{signing.py}: \texttt{sign}, \texttt{verify} and the keys of Alice and Mallory in PEM files.}
  \vfill
  \opentask{Task2}{about 10 minutes}
  \vfill
  \say{Około 10 minut. Polecenia w README. Najpierw sprawdzić, czy wszyscy mają bibliotekę: python3 -c "import cryptography". Pięć TODO: podpis, sprawdzenie, zmieniona wiadomość (11:00 zamiast 10:00), podpis Mallory jej kluczem sprawdzany kluczem Alice, drugi podpis tej samej wiadomości. Przed TODO 5 zapytać: czy drugi podpis będzie taki sam? Ed25519 jest deterministyczny, więc wszyscy w sali dostaną bajt w bajt ten sam podpis: można porównać pierwsze znaki z sąsiadem.}
\end{frame}

\samedot
\begin{frame}{Task 2: what we saw}
  % CHECKED: python3 Task2/solution/sign_and_verify.py, 2026-11-17
  \vfill
  \stepbox{The signature, 64 bytes}{\ttfamily\scriptsize eb232a684f057336cd2780944693c1992e27bda5f6a383416f4a1caa2739a8b5\\4298b02954b5a4d43f647beaf3f1163e3fdea3aa7241cc9a95fcdf85a69d7f06}
  \vfill
  \begin{center}
    \renewcommand{\arraystretch}{1.4}
    \begin{tabular}{@{}l c@{}}
      \toprule
      Bob checks the message & \ok\ \texttt{True} \\
      Bob checks the message with 11:00 instead of 10:00 & \nok\ \texttt{False} \\
      Bob checks Mallory's signature with Alice's public key & \nok\ \texttt{False} \\
      the second signature of the same message is the same & \texttt{True} \\
      \bottomrule
    \end{tabular}
  \end{center}
  \vfill
  \keymsg{One changed character, or a signature from another key: verify says False.}
  \say{Wynik rozwiązania. Wszyscy powinni mieć ten sam podpis eb232a68... (Ed25519 jest deterministyczny, klucze z plików są stałe). Mallory umie podpisać, ale swoim kluczem; sprawdzenie kluczem Alice tego nie przepuści. Zapytać: kto mógł zrobić podpis, który przechodzi z kluczem publicznym Alice? (Tylko ktoś z kluczem prywatnym Alice.)}
\end{frame}

\begin{frame}{Who wrote it: a first answer}
  \vfill
  \begin{center}
    \begin{tikzpicture}[>=Stealth]
      \alice{0,0}
      \node[apriv, font=\small] at (0,-1.6) {\faKey\ private key};
      \node[asig, font=\small] (s) at (4.2,0) {\faStamp\ signature};
      \draw[data] (A.east) to node[above, font=\scriptsize] {only she can make it} (s.west);
      \foreach \n/\x in {Bob/8.0, Carol/9.6, Dave/11.2} {
        \node[circle, fill=nLine, text=white, minimum size=0.9cm, font=\normalsize] (p\n) at (\x,0) {\faUser};
        \node[font=\scriptsize\bfseries, text=nInk, below=1pt of p\n] {\n};
        \node[apub, font=\tiny, inner sep=2pt] at (\x,-1.3) {\faKey\ public};
      }
      \draw[data] (s.east) to node[above, font=\scriptsize] {everybody checks} (7.4,0);
    \end{tikzpicture}
  \end{center}
  \vfill
  \bigdef{A message with a valid signature was written, or at least approved, by \textbf{whoever holds the private key}.
    Not Mallory, not Bob, not Carol. And anybody can check it, without any secret.}
  \vfill
  \say{Pierwsza odpowiedź na główne pytanie wykładu. Podpis rozwiązuje oba problemy MAC: Bob nie umie podrobić podpisu Alice (nie ma jej klucza prywatnego), a sprawdzać mogą wszyscy, także sędzia Carol. ,,Or at least approved'': podpis mówi, że właściciel klucza zaakceptował te bajty, a nie, że je sam wymyślił. Dokładniejszą odpowiedź, z zastrzeżeniem o kluczu publicznym, damy na końcu. Teraz zajrzymy do środka sign i verify.}
\end{frame}

% ======================================================================
\section[RSA signatures inside]{RSA signatures inside}
% ======================================================================

\begin{frame}{Three ways to sign}
  \vfill
  \begin{center}
    \begin{tikzpicture}
      \foreach \x/\name/\yr/\base [count=\k] in {0/RSA/1977/{factoring\\(Lecture 3)}, 4.6/ECDSA/1992/{points on a curve\\(Lecture 5)}, 9.2/Ed25519/2011/{points on a curve\\(Lecture 5)}} {
        \node[hw, minimum width=4.0cm, minimum height=2.6cm] (b\k) at (\x,0) {};
        \node[font=\large\bfseries, text=nInk] at (\x,0.75) {\name};
        \node[font=\small, text=gray] at (\x,0.25) {\yr};
        \node[font=\small, text=nInk, align=center] at (\x,-0.55) {\base};
      }
      \draw[annot] (-2.1,-1.4) rectangle (2.1,1.4);
      \node[font=\small\bfseries, text=AccentOrange] at (0,-1.75) {we open this one};
    \end{tikzpicture}
  \end{center}
  \vfill
  \begin{center}\small All three offer the same sign and verify. RSA is the simplest to open: we already know its keys.\end{center}
  \keymsg{Inside RSA signatures: the same two powers as RSA encryption.}
  % CHECKED: RSA: Rivest, Shamir, Adleman, MIT memo 1977, CACM 1978; ECDSA: proposed by Vanstone in 1992 (FIPS 186-5 history), Ed25519: Bernstein et al. 2011
  \say{RSA podpisuje od początku: artykuł Rivesta, Shamira i Adlemana z 1977/1978 miał w tytule ,,digital signatures''. ECDSA zaproponował Scott Vanstone w 1992 roku, Ed25519 to rok 2011. ECDSA i Ed25519 zostaną w sekcji 3 czarnymi skrzynkami, poza jednym atakiem.}
\end{frame}

\begin{frame}{Recall: RSA}
  \vfill
  \begin{center}
    \begin{tikzpicture}[>=Stealth]
      \node[bpub, font=\Large, minimum width=4.2cm, minimum height=1.4cm] at (0,0) {\faLockOpen\ \{e, n\} = \{7, 187\}\\[2pt]\small public};
      \node[bpriv, font=\Large, minimum width=4.2cm, minimum height=1.4cm] at (5.6,0) {\faKey\ \{d, n\} = \{23, 187\}\\[2pt]\small private};
      \node[font=\large, text=nInk] at (0,-1.5) {encrypt: $C = P^{e} \bmod n$};
      \node[font=\large, text=nInk] at (5.6,-1.5) {decrypt: $P = C^{d} \bmod n$};
    \end{tikzpicture}
  \end{center}
  \vfill
  \begin{center}\small
    $n = 17 \times 11 = 187$, \quad $e \times d = 7 \times 23 = 161 = 160 + 1$:\\
    one full circle of powers plus one step, so $(P^{e})^{d} = P$.
  \end{center}
  \keymsg{Two powers modulo n that undo each other.}
  \say{Przypomnienie z wykładu 3, ten sam przykład z Stallingsa: p = 17, q = 11, phi = 160, e = 7, d = 23. Kluczowy fakt: e razy d = 161, czyli 1 więcej niż phi, więc potęga do e, a potem do d daje z powrotem P. Zapytać: czy kolejność potęg ma znaczenie? (Nie: e razy d = d razy e.)}
\end{frame}

\samedot
\begin{frame}{The RSA signature}
  \vfill
  \begin{center}
    \begin{tikzpicture}[>=Stealth]
      \node[msg, font=\LARGE, minimum width=1.6cm] (m) at (0,0) {M};
      \node[asig, font=\LARGE, minimum width=1.6cm] (s) at (6,0) {S};
      \node[msg, font=\LARGE, minimum width=1.6cm] (m2) at (12,0) {M?};
      \draw[data] (m) to node[above=6pt, apriv, font=\normalsize] {\faKey\ $S = M^{d} \bmod n$} (s);
      \draw[data] (s) to node[above=6pt, apub, font=\normalsize] {\faKey\ $S^{e} \bmod n$} (m2);
      \node[font=\scriptsize, text=nInk] at (3,-0.6) {sign with the private key};
      \node[font=\scriptsize, text=nInk] at (9,-0.6) {verify with the public key};
    \end{tikzpicture}
  \end{center}
  \vfill
  \bigdef{\textbf{Sign}: $S = M^{d} \bmod n$. \qquad \textbf{Verify}: compute $S^{e} \bmod n$; \texttt{True} if it equals $M$.}
  \vfill
  \begin{center}\small Now the key belongs to Alice, the signer: the same numbers, Alice's colour.\end{center}
  \keymsg{Signing is the private power; verifying is the public power and a comparison.}
  \say{Te same dwie potęgi co przy szyfrowaniu, tylko w odwrotnej kolejności: najpierw prywatna, potem publiczna. M to wiadomość jako liczba mniejsza od n (jak P w wykładzie 3). Bob dostaje M i S, liczy S do e mod n i porównuje z M. Klucz z wykładu 3 należał do Boba; dziś podpisuje Alice, więc te same liczby są w jej kolorze. Uwaga na mit: ,,podpis to szyfrowanie kluczem prywatnym'' działa tylko dla podręcznikowego RSA; dla ECDSA i Ed25519 nie ma takiej analogii, więc lepiej mówić sign i verify.}
\end{frame}

\samedot
\begin{frame}{Why the signature checks}
  \vfill
  \stepbox{1. What Bob computes}{$S^{e} = (M^{d})^{e} = M^{d \times e} \pmod n$}
  \stepbox{2. The keys were made for this (Lecture 3)}{$d \times e = 23 \times 7 = 161 = 160 + 1$: one full circle of powers, plus one step.}
  \stepbox{3. So the message comes back}{$M^{161} = M^{160} \times M = 1 \times M = M \pmod{187}$}
  \vfill
  \begin{columns}[t,onlytextwidth]
    \column{0.48\textwidth}\plainbox{\ok\ Alice: she knows $d$, so she can make $S$.}
    \column{0.48\textwidth}\plainbox{\nok\ Mallory: to make $S$ without $d$ she would have to undo a power modulo $n$, as hard as breaking RSA.}
  \end{columns}
  \keymsg{It works for the same reason RSA decryption works.}
  \say{Dowód dokładnie jak na slajdzie Why does the message come back z wykładu 3, tylko z zamienioną kolejnością wykładników. Dlaczego Mallory nie podrobi: musiałaby znaleźć S takie, że S do 7 mod n = M, czyli pierwiastek siódmego stopnia modulo n. Bez d (czyli bez rozkładu n) nikt nie zna szybkiej metody. Zaraz zobaczymy jednak, że to ,,dla wybranego M''; Mallory ma sprytniejsze sposoby.}
\end{frame}

\samedot
\begin{frame}{The example: 11 $\to$ 88 $\to$ 11}
  \vfill
  \begin{center}
    \begin{tikzpicture}[>=Stealth]
      \node[msg, font=\LARGE] (m) at (0,0) {11};
      \node[asig, font=\LARGE] (s) at (5.5,0) {88};
      \node[msg, font=\LARGE] (m2) at (11,0) {11};
      \draw[data] (m) to node[above=6pt, apriv, font=\small] {\faKey\ \{23, 187\}} (s);
      \draw[data] (s) to node[above=6pt, apub, font=\small] {\faKey\ \{7, 187\}} (m2);
      \node[font=\normalsize, text=nInk] at (2.75,-0.9) {$S = 11^{23} \bmod 187 = \mathbf{88}$};
      \node[font=\normalsize, text=nInk] at (8.25,-0.9) {$88^{7} \bmod 187 = \mathbf{11}$ \;\ok};
      \node[font=\small\ttfamily, text=nInk] at (2.75,-1.6) {pow(11, 23, 187)};
      \node[font=\small\ttfamily, text=nInk] at (8.25,-1.6) {pow(88, 7, 187)};
    \end{tikzpicture}
  \end{center}
  \vfill
  \begin{center}
    Mallory changes the message to 12 and keeps $S = 88$:\quad $88^{7} \bmod 187 = 11 \neq 12$ \;\nok
  \end{center}
  \vfill
  \curious{signatures/\#2}
  \keymsg{The signature fits one message; change the message and verify fails.}
  % CHECKED: python3: pow(11, 23, 187) = 88, pow(88, 7, 187) = 11
  \say{Liczby sprawdzone w Pythonie. Ciekawostka: to są liczby z wykładu 3 w odwrotną stronę (tam 88 szyfrowało się do 11). Bob z podpisem 88 i zmienioną wiadomością 12 dostaje 11, więc odrzuca. W symulatorze (link na końcu sekcji) można przeklikać wszystkie M od 0 do 186.}
\end{frame}

\samedot
\begin{frame}{Textbook RSA signatures}
  \vfill
  \bigdef{This bare formula, $S = M^{d} \bmod n$ with nothing added, is called a \textbf{textbook RSA signature}.}
  \vfill
  \begin{columns}[t,onlytextwidth]
    \column{0.48\textwidth}\stepbox{A message as a number}{as in Lecture 3: the bytes read as one big-endian number, smaller than $n$.\\ \texttt{"Hi"} $\to$ \texttt{0x4869} = 18\,537}
    \column{0.48\textwidth}\stepbox{A real key}{$n$ has 2048 bits (256 bytes): $d$, $S$ and $n$ are numbers with about 617 digits. \texttt{pow} handles them in a millisecond.}
  \end{columns}
  \vfill
  \keymsg{Textbook RSA works; in the next section we will see why nobody uses it alone.}
  % CHECKED: python3: int.from_bytes(b"Hi", "big") = 18537; 2**2048 has 617 digits
  \say{Nazwa jak przy szyfrowaniu w wykładzie 3 (textbook RSA). Zadanie 3 robi to na prawdziwym kluczu 2048 bitów Alice: liczby mają 617 cyfr dziesiętnych (512 cyfr szesnastkowych), Python liczy na nich bez problemu. Po zadaniu zobaczymy trzy problemy podręcznikowego podpisu.}
\end{frame}

\begin{frame}{Task 3: textbook RSA}
  \vfill
  \taskgoal{sign ``Pay Bob 10 pounds. Alice'' with \texttt{pow}, check it, then check a changed copy.}
  \taskready{\texttt{alice\_key.py}: Alice's 2048-bit key as the numbers \texttt{n}, \texttt{e}, \texttt{d}.}
  \vfill
  \opentask{Task3}{about 8 minutes}
  \vfill
  \say{Około 8 minut. Pięć krótkich TODO: wiadomość jako liczba (int.from\_bytes), podpis pow(m, d, n), sprawdzenie pow(s, e, n) == m, zmieniona wiadomość (90 zamiast 10 funtów), a w TODO 5 pytanie: co jest w środku podpisu? Przed TODO 5 zapytać salę, co wyjdzie z pow(s, e, n) zamienionego na bajty. Kto skończy wcześniej: niech zmieni wiadomość na dłuższą niż 256 bajtów i zobaczy, co się stanie (m nie jest mniejsze od n).}
\end{frame}

\samedot
\begin{frame}{Task 3: what we saw}
  % CHECKED: python3 Task3/solution/textbook_rsa.py, 2026-11-17
  \vfill
  \begin{center}
    \renewcommand{\arraystretch}{1.4}
    \begin{tabular}{@{}l l@{}}
      \toprule
      $m$ (the message as a number) & \texttt{50617920426f\ldots 20416c696365}, 48 hex digits \\
      $s = m^{d} \bmod n$ & \texttt{1ee3425e9a74\ldots 95e50c23a9ae}, 512 hex digits \\
      Bob checks the message & \ok\ \texttt{True} \\
      Bob checks ``\ldots 90 pounds \ldots'' with the same $s$ & \nok\ \texttt{False} \\
      $s^{e} \bmod n$ as bytes & \texttt{b'Pay Bob 10 pounds. Alice'} \\
      \bottomrule
    \end{tabular}
  \end{center}
  \vfill
  \begin{center}\small The signature carries the message: $s^{e} \bmod n$ gives back exactly what Alice signed.\end{center}
  \keymsg{The signature is as long as $n$ (256 bytes), even for a 24-byte message.}
  \say{Wynik rozwiązania. 50617920 to bajty P, a, y, spacja w zapisie szesnastkowym: wiadomość jako liczba. Podpis ma 512 cyfr szesnastkowych, czyli 256 bajtów, tyle co n. TODO 5: z podpisu i klucza publicznego wychodzi oryginalna wiadomość, więc Bob widzi, co Alice naprawdę podpisała (10 funtów, nie 90). Klucz w pliku jest ten sam dla wszystkich, więc wszyscy mają ten sam podpis.}
\end{frame}

\begin{frame}{Problem 1: long messages}
  % CHECKED: python3: ceil(10 000 000 / 255) = 39 216 pieces, x 256 bytes = 10 039 296 bytes; x 1.1 ms (Task 7, RSA-2048 sign) = 43 s
  \vfill
  \begin{center}
    \begin{tikzpicture}
      \node[msg, minimum width=11.5cm, minimum height=0.8cm, font=\small] (f) at (0,0) {\faFile*[regular]\ an update of 10 MB};
      \foreach \i in {0,...,9} \node[msg, minimum width=0.95cm, minimum height=0.6cm, font=\tiny] at ({-5.3+1.05*\i},-1.3) {255 B};
      \node[font=\small, text=nInk] at (5.4,-1.3) {\ldots};
      \foreach \i in {0,...,9} \node[asig, minimum width=0.95cm, minimum height=0.6cm, font=\tiny] at ({-5.3+1.05*\i},-2.3) {256 B};
      \node[font=\small, text=nInk] at (5.4,-2.3) {\ldots};
      \node[font=\scriptsize, text=gray, anchor=west] at (6.0,-1.3) {39\,216 pieces};
      \node[font=\scriptsize, text=gray, anchor=west] at (6.0,-2.3) {39\,216 signatures};
    \end{tikzpicture}
  \end{center}
  \vfill
  \begin{itemize}\small
    \item $M$ must be smaller than $n$: at most 255 bytes per signature with a 2048-bit key.
    \item 10 MB need \textbf{39\,216 signatures}: 10 MB of signatures and about 43 seconds of signing.
    \item Each piece is signed alone: Mallory can \textbf{drop or reorder} pieces, and every signature still checks.
  \end{itemize}
  \keymsg{Textbook RSA signs only short numbers.}
  \say{10 000 000 / 255 = 39 216 kawałków (w górę). Każdy podpis ma 256 bajtów, razem około 10 MB, tyle samo co plik. Czas: 1,1 ms na podpis RSA-2048 (pomiar z zadania 7 na i5-8265U) razy 39 216 = około 43 s. A najgorsze: kawałki są podpisane osobno, więc Mallory może wyrzucić kawałek ze środka albo zamienić kolejność, i każdy podpis nadal jest poprawny.}
\end{frame}

\samedot
\begin{frame}{Problem 2: signatures multiply}
  % CHECKED: python3: pow(2, 23, 187) = 162, pow(3, 23, 187) = 181, 162 * 181 % 187 = 150 = pow(6, 23, 187), pow(150, 7, 187) = 6
  \vfill
  \begin{center}
    \begin{tikzpicture}[>=Stealth]
      \node[msg, font=\Large] (m2) at (0,1) {2};
      \node[asig, font=\Large] (s2) at (2.6,1) {162};
      \node[msg, font=\Large] (m3) at (0,-0.4) {3};
      \node[asig, font=\Large] (s3) at (2.6,-0.4) {181};
      \draw[data] (m2) -- (s2); \draw[data] (m3) -- (s3);
      \node[font=\scriptsize, text=gray] at (1.3,1.9) {signed by Alice};
      \node[msg, font=\Large] (m6) at (9.4,0.3) {6};
      \node[mpub, font=\Large] (s6) at (6.6,0.3) {150};
      \draw[data] (s2.east) -- (5.9,0.45);
      \draw[data] (s3.east) -- (5.9,0.15);
      \node[font=\scriptsize, text=sPinkText, align=center] at (6.6,1.4) {Mallory: $162 \times 181 \bmod 187$};
      \draw[data] (s6) to node[above, font=\scriptsize] {$150^{7} \bmod 187$} (m6);
      \node[font=\Large] at (10.3,0.3) {\ok};
    \end{tikzpicture}
  \end{center}
  \vfill
  \bigdef{$M_1^{d} \times M_2^{d} = (M_1 \times M_2)^{d} \pmod n$: \;the product of two signatures is a signature of the product.}
  \vfill
  \keymsg{From two signed messages, Mallory makes a third signature without the private key.}
  \say{Rachunek: S(2) = 2 do 23 mod 187 = 162, S(3) = 181. Mallory mnoży: 162 razy 181 mod 187 = 150. I rzeczywiście 150 do 7 mod 187 = 6, więc 150 to poprawny podpis wiadomości 6, której Alice nigdy nie podpisała. Dlaczego: potęgowanie zachowuje mnożenie, (ab) do d = a do d razy b do d. Podpis Mallory jest różowy: zrobiła go ona, choć przechodzi z kluczem Alice.}
\end{frame}

\samedot
\begin{frame}{Problem 3: any number is a signature}
  % CHECKED: python3: pow(100, 7, 187) = 144, pow(144, 23, 187) = 100
  \vfill
  \begin{center}
    \begin{tikzpicture}[>=Stealth]
      \mallory{0,0}
      \node[mpub, font=\LARGE] (s) at (3.4,0) {100};
      \node[font=\scriptsize, text=gray] at (3.4,0.75) {any number as $S$};
      \node[msg, font=\LARGE] (m) at (8.6,0) {144};
      \draw[data] (M.east) -- (s.west);
      \draw[data] (s) to node[above, apub, font=\small] {\faKey\ $100^{7} \bmod 187$} (m);
      \node[font=\scriptsize, text=gray] at (8.6,0.75) {the message it signs};
      \node[font=\large, anchor=west] at (9.4,0) {\ok};
    \end{tikzpicture}
  \end{center}
  \vfill
  \begin{center}
    Mallory goes backwards: she picks $S$ first, then computes $M = S^{e} \bmod n$ with the \textbf{public} key.\\[4pt]
    The pair $(M, S) = (144, 100)$ verifies. She cannot choose $M$, but\\
    if any number is a fine message (a random code, an amount), that is enough.
  \end{center}
  \keymsg{Textbook RSA lets anybody make valid pairs; only the choice of the message is hard.}
  \say{Trzeci problem: Mallory nie wybiera wiadomości, tylko podpis. 100 do 7 mod 187 = 144, i para (144, 100) przechodzi sprawdzenie. Jeśli system przyjmuje dowolną liczbę (np. jednorazowy kod, losowy identyfikator, kwotę), to Mallory ma poprawnie podpisaną liczbę. Zadanie 4 robi oba triki na prawdziwym kluczu 2048 bitów.}
\end{frame}

\begin{frame}{Task 4: forge it}
  \vfill
  \taskgoal{as Mallory, without the private key, make Alice's bank accept an order for 1000 pounds.}
  \taskready{\texttt{bank\_ready.py}: Alice's public key and two orders she signed for Mallory, 20 and 50 pounds.}
  \vfill
  \opentask{Task4}{about 9 minutes}
  \vfill
  \say{Około 9 minut. Bank Alice przyjmuje zlecenie jako liczbę (kwota w funtach) podpisaną podręcznikowym RSA. TODO 1: bank sprawdza dwa prawdziwe zlecenia. TODO 2: mnożenie podpisów, 20 razy 50 = 1000. TODO 3: dowolna liczba jako podpis. Jeśli ktoś utknie przy TODO 2: przypomnieć slajd Problem 2, ale na prawdziwym kluczu.}
\end{frame}

\samedot
\begin{frame}{Task 4: what we saw}
  % CHECKED: python3 Task4/solution/forge.py, 2026-11-17
  \vfill
  \begin{center}
    \renewcommand{\arraystretch}{1.4}
    \begin{tabular}{@{}l l c@{}}
      \toprule
      order 1, 20 pounds & signed by Alice & \ok\ \texttt{True} \\
      order 2, 50 pounds & signed by Alice & \ok\ \texttt{True} \\
      $s_1 \times s_2 \bmod n$ (\texttt{13bf50ca\ldots 1031c3}) & on \textbf{1000 pounds} & \ok\ \texttt{True} \\
      any $s$: 123\,456\,789 & on \texttt{4383e8aa\ldots 62c032b0} & \ok\ \texttt{True} \\
      \bottomrule
    \end{tabular}
  \end{center}
  \vfill
  \begin{center}\small Mallory never had $d$. Two signatures and one multiplication were enough.\end{center}
  \keymsg{Textbook RSA signatures can be forged: we must not sign the message itself.}
  \say{Wynik rozwiązania. Podpis na 1000 funtów: iloczyn podpisów mod n, bank mówi True. Losowe S = 123456789 podpisuje ogromną liczbę (512 cyfr szesnastkowych); bank też mówi True, choć kwota jest bez sensu. Zapytać: jak to naprawić? Podpowiedź: na początku wykładu mieliśmy coś, co z każdego pliku robi 32 bajty i czego nie da się odwrócić.}
\end{frame}

\begin{frame}{Hash first, then sign}
  \vfill
  \begin{center}
    \begin{tikzpicture}[>=Stealth]
      \node[msg, font=\small, minimum width=2.4cm, minimum height=1.6cm] (m) at (0,0) {\faFile*[regular]\\ message\\ any length};
      \mech{h}{3.2,0}{SHA-256}{2.2cm}
      \node[digest, font=\small, minimum width=2.0cm] (d) at (6.0,0) {digest\\32 bytes};
      \node[hwmain, minimum width=2.4cm, minimum height=1.0cm, font=\small\bfseries] (r) at (9.0,0) {\faCogs\ $H^{d} \bmod n$};
      \node[apriv, font=\small] (k) at (9.0,1.3) {\faKey\ private key};
      \node[asig, font=\small, minimum height=0.9cm] (s) at (12.0,0) {\faStamp\ signature};
      \draw[data] (m) -- (h); \draw[data] (h) -- (d); \draw[data] (d) -- (r); \draw[data] (k) -- (r); \draw[data] (r) -- (s);
    \end{tikzpicture}
  \end{center}
  \vfill
  \bigdef{\textbf{Hash-and-sign}: \;$S = H(M)^{d} \bmod n$. \quad Verify: $S^{e} \bmod n = H(M)$?\\[2pt]
    \small Bob computes $H(M)$ himself from the message he received.}
  \vfill
  \keymsg{We sign the digest of the message, not the message.}
  \say{Poprawka: zamiast M podpisujemy H(M), skrót wiadomości (SHA-256), czytany jako liczba. Bob dostaje wiadomość i podpis, liczy skrót sam i porównuje z S do e mod n. Tak robi każdy prawdziwy podpis, także ECDSA i Ed25519 (Ed25519 używa SHA-512 w środku). Dlatego na wykładzie 9 było napisane, że w środku liczy się skrót danych. Brązowy to skrót: ten sam kolor co na początku wykładu.}
\end{frame}

\samedot
\begin{frame}{Bob checks a hash-and-sign signature}
  \vfill
  \begin{center}
    \begin{tikzpicture}[>=Stealth]
      \bob{-1.6,-0.6}
      \node[msg, font=\small] (m) at (1.4,0.6) {\faFile*[regular]\ message he got};
      \mech{h}{5.0,0.6}{SHA-256}{2.2cm}
      \node[digest, font=\small] (d1) at (8.6,0.6) {$H(M)$};
      \node[asig, font=\small] (s) at (1.4,-1.0) {\faStamp\ signature $S$};
      \node[hwmain, minimum width=2.2cm, minimum height=1.0cm, font=\small\bfseries] (p) at (5.0,-1.0) {\faCogs\ $S^{e} \bmod n$};
      \node[apub, font=\scriptsize] at (5.0,-1.95) {\faKey\ Alice's public key};
      \node[digest, font=\small] (d2) at (8.6,-1.0) {a number};
      \draw[data] (m) -- (h); \draw[data] (h) -- (d1);
      \draw[data] (s) -- (p); \draw[data] (p) -- (d2);
      \node[font=\Large, anchor=west] at (9.8,-0.2) {equal? \ok\,/\,\nok};
    \end{tikzpicture}
  \end{center}
  \vfill
  \begin{center}\small
    Two independent ways to the same 32 bytes: from the message (anybody), and from the signature (only Alice could start it).\\
    If they meet, the message is the one Alice signed.
  \end{center}
  \keymsg{Verify = hash the message yourself, undo the signature with the public key, compare.}
  \say{Strona Boba. Bob nie dostaje skrótu od nikogo: liczy go sam z wiadomości, którą otrzymał. Z podpisu, potęgą publiczną, wychodzi liczba, którą Alice podpisała. Jeśli Mallory zmieniła wiadomość, skrót Boba będzie inny; jeśli zmieniła podpis, wyjdzie inna liczba. W zadaniu 5 TODO 3 to dokładnie ten rysunek.}
\end{frame}

\samedot
\begin{frame}{Why the hash helps: any length}
  \vfill
  \begin{center}
    \begin{tikzpicture}[>=Stealth]
      \node[msg, minimum width=8cm, minimum height=0.8cm, font=\small] (f) at (0,0) {\faFile*[regular]\ an update of 10 MB};
      \mech{h}{0,-1.4}{SHA-256}{2.2cm}
      \node[digest, font=\small] (d) at (0,-2.6) {32 bytes};
      \node[asig, font=\small] (s) at (4.0,-2.6) {\faStamp\ one signature, 256 bytes};
      \draw[data] (f) -- (h); \draw[data] (h) -- (d); \draw[data] (d) -- (s);
    \end{tikzpicture}
  \end{center}
  \vfill
  \begin{columns}[t,onlytextwidth]
    \column{0.32\textwidth}\stepbox{One signature}{instead of 39\,216: one power, about 1 ms.}
    \column{0.32\textwidth}\stepbox{The whole file}{change, drop or move one byte and the digest is completely different.}
    \column{0.32\textwidth}\stepbox{Always small}{a 32-byte digest is always smaller than $n$.}
  \end{columns}
  \keymsg{One digest covers the whole file, so one signature does too.}
  \say{Problem 1 rozwiązany: skrót ma zawsze 32 bajty, niezależnie od rozmiaru pliku. Skrót 10 MB liczy się w ułamku sekundy (SHA-256 robi kilkaset MB na sekundę), a potęga jest jedna. Nie da się też usuwać ani przestawiać kawałków, bo podpis obejmuje skrót całości.}
\end{frame}

\samedot
\begin{frame}{Why the hash helps: no forgeries}
  \vfill
  \begin{center}
    \begin{tikzpicture}[>=Stealth]
      \node[asig, font=\small] (s1) at (0,0.6) {$S_1$ on $H(M_1)$};
      \node[asig, font=\small] (s2) at (0,-0.6) {$S_2$ on $H(M_2)$};
      \node[mpub, font=\small] (s) at (4.0,0) {$S_1 \times S_2$};
      \draw[data] (s1.east) -- (s.west); \draw[data] (s2.east) -- (s.west);
      \node[digest, font=\small] (n) at (8.0,0) {the number $H(M_1) \times H(M_2) \bmod n$};
      \draw[data] (s) to node[above, font=\scriptsize] {signs} (n);
      \node[msg, font=\small, dashed] (q) at (8.0,-1.6) {a message $M$ with this digest?};
      \draw[data, dashed] (n) -- (q);
      \node[font=\small, text=aWrite, anchor=west] at (10.3,-1.6) {\nok\ one way};
    \end{tikzpicture}
  \end{center}
  \vfill
  \begin{center}\small
    The multiplied or random signature still signs some \textbf{number}. To use it, Mallory needs a message\\
    whose SHA-256 is exactly that number: she must run SHA-256 backwards (Lecture 6: one way).
  \end{center}
  \keymsg{After hashing, a forged number is useless: no message leads to it.}
  \say{Problemy 2 i 3. Mnożenie nadal działa na liczbach, ale wynik podpisuje liczbę, która nie jest skrótem żadnej znanej wiadomości. Żeby ją wykorzystać, Mallory musiałaby znaleźć wiadomość o takim skrócie, czyli odwrócić SHA-256 (preimage, wykład 6): około $2^{256}$ prób. Losowe S tak samo. Zadanie 5 to sprawdza.}
\end{frame}

\begin{frame}{Task 5: hash, then sign}
  \vfill
  \taskgoal{sign a 100\,000-byte update with hash-and-sign; change one byte; try the multiplication trick again.}
  \taskready{\texttt{alice\_key.py} (\texttt{n}, \texttt{e}, \texttt{d}); \texttt{update.bin}; two hash-signed orders in \texttt{signed\_orders.txt}.}
  \vfill
  \opentask{Task5}{about 9 minutes}
  \vfill
  \say{Około 9 minut. TODO 1: funkcja hash\_number (SHA-256 jako liczba). TODO 2 i 3: podpis i sprawdzenie. TODO 4: zmiana jednego bitu w bajcie 50 000 (bytearray, ^= 1). TODO 5: ten sam trik z mnożeniem co w zadaniu 4, tym razem na podpisach skrótów. Przed TODO 5 zapytać: czy bank przyjmie 1000 funtów?}
\end{frame}

\samedot
\begin{frame}{Task 5: what we saw}
  % CHECKED: python3 Task5/solution/hash_and_sign.py, 2026-11-17
  \vfill
  \begin{center}
    \renewcommand{\arraystretch}{1.4}
    \begin{tabular}{@{}l l@{}}
      \toprule
      the file, 100\,000 bytes & digest \texttt{710fcd1df5b9\ldots 3416108c9faf} \\
      one signature, 256 bytes & \texttt{bf0aa615a40e\ldots f8609c05042f} \\
      Bob checks the file & \ok\ \texttt{True} \\
      Bob checks it with one bit of byte 50\,000 flipped & \nok\ \texttt{False} \\
      the bank checks $s_{20} \times s_{50}$ on ``\ldots 1000 pounds'' & \nok\ \texttt{False} \\
      \bottomrule
    \end{tabular}
  \end{center}
  \vfill
  \keymsg{Hash-and-sign: one signature for any file, and the multiplication trick is dead.}
  \say{Wynik rozwiązania. Plik 100 000 bajtów, jeden podpis 256 bajtów. Jeden bit zmieniony: skrót zupełnie inny, więc False. Iloczyn podpisów nie pasuje do skrótu wiadomości o 1000 funtów, więc bank odrzuca. Zapytać: czy to znaczy, że podręcznikowe RSA na skrócie jest już idealne? Odpowiedź na następnym slajdzie.}
\end{frame}

\begin{frame}{In practice: RSA-PSS}
  \vfill
  \begin{center}
    \begin{tikzpicture}[>=Stealth]
      \node[digest, font=\small, minimum width=2.2cm] (d) at (0,0) {digest};
      \node[font=\large] at (1.55,0) {+};
      \node[hw, font=\small, minimum width=1.8cm] (s) at (3.0,0) {\faDice\ salt};
      \node[font=\large] at (4.45,0) {+};
      \node[hw, font=\small, minimum width=2.4cm] (p) at (6.1,0) {fixed structure};
      \node[hwmain, font=\small\bfseries] (r) at (9.4,0) {\faCogs\ power $d$};
      \draw[data] (p) -- (r);
      \node[asig, font=\small] (sig) at (12.0,0) {\faStamp\ signature};
      \draw[data] (r) -- (sig);
    \end{tikzpicture}
  \end{center}
  \vfill
  \bigdef{\textbf{RSA-PSS} does for signatures what OAEP does for encryption (Lecture 3):\\
    random salt and a fixed structure around the digest before the power, so no number games work.}
  \vfill
  \begin{center}\small
    The older \textbf{PKCS\#1 v1.5} padding (fixed, no salt) is still in most certificates today.
    Never textbook RSA.
  \end{center}
  \keymsg{Real RSA signatures: hash, pad, then power.}
  % CHECKED: RFC 8017 (PKCS #1 v2.2), section 8.1 RSASSA-PSS and 8.2 RSASSA-PKCS1-v1_5;
  %   RFC 8446 section 4.2.3: rsa_pkcs1_* "refer solely to signatures which appear in certificates" (bugzilla.mozilla.org 1489997, 2026-11-17)
  \say{Jedno zdanie, jak OAEP w wykładzie 3: PSS (Bellare i Rogaway, 1996) dodaje losową sól i stałą strukturę bajtów, zanim liczba pójdzie do potęgi. Dzięki soli ten sam plik podpisany dwa razy daje różne podpisy (zobaczymy w zadaniu 7). Starsze dopełnienie PKCS\#1 v1.5 jest deterministyczne i nadal bezpieczne do podpisów, i jest w większości certyfikatów. W TLS 1.3 (RFC 8446, punkt 4.2.3) serwer, który podpisuje RSA w trakcie łączenia, musi używać PSS; v1.5 zostało tylko dla podpisów w certyfikatach. W bibliotece cryptography: padding.PSS.}
\end{frame}

\samedot
\begin{frame}[fragile]{RSA-PSS in Python}
\begin{lstlisting}[style=py]
from cryptography.hazmat.primitives import hashes
from cryptography.hazmat.primitives.asymmetric import padding

pss = padding.PSS(mgf=padding.MGF1(hashes.SHA256()), salt_length=32)

signature = private_key.sign(data, pss, hashes.SHA256())   # 256 bytes
public_key.verify(signature, data, pss, hashes.SHA256())   # exception if wrong
\end{lstlisting}
  \vfill
  \begin{center}\small
    The library hashes the data, adds the salt and the structure, and computes the power:\\
    everything of this section in one call. Signed twice, the same data gives two different signatures.
  \end{center}
  \keymsg{In real code: the library and RSA-PSS, never your own \texttt{pow}.}
  \say{To samo co pow z zadań 3 do 5, ale poprawnie: skrót, sól, struktura, potęga. MGF1 to sposób rozciągnięcia skrótu na długość n (nie trzeba pamiętać). Zadanie 7 używa właśnie tego podpisu jako RSA-2048. Pow w zadaniach był po to, żeby zobaczyć środek; w prawdziwym programie nigdy nie liczymy podpisu sami.}
\end{frame}

\begin{frame}{Recall: a collision}
  \vfill
  \begin{center}
    \begin{tikzpicture}[>=Stealth]
      \node[msg, font=\small, minimum width=2.8cm] (a) at (0,0.7) {document A};
      \node[msg, font=\small, minimum width=2.8cm] (b) at (0,-0.7) {document B};
      \mech{h}{3.8,0}{MD5}{1.8cm}
      \node[digest, font=\small] (d) at (7.2,0) {the same digest};
      \draw[data] (a.east) -- (h.west); \draw[data] (b.east) -- (h.west); \draw[data] (h) -- (d);
    \end{tikzpicture}
  \end{center}
  \vfill
  \begin{center}\small
    Lecture 7: for MD5 and SHA-1, anyone can make two different inputs with the same digest,\\
    even with a chosen beginning (chosen prefix). SHA-256 has no known collision.
  \end{center}
  \keymsg{A broken hash gives two inputs with one digest.}
  \say{Przypomnienie z wykładu 7: kolizje MD5 (Wang 2004) i SHA-1 (SHAttered 2017, Shambles 2020, z wybranym prefiksem). SHA-256 trzyma się: najlepszy znany atak to nadal około $2^{128}$ prób (paradoks urodzin). Co to znaczy dla podpisu, który podpisuje tylko skrót?}
\end{frame}

\samedot
\begin{frame}{A collision breaks the signature}
  \vfill
  \begin{center}
    \begin{tikzpicture}[>=Stealth]
      \node[msg, font=\small, text width=3.2cm] (a) at (0,0.9) {A: ``a website certificate''};
      \node[msg, font=\small, text width=3.2cm, draw=sPink, line width=1.2pt] (b) at (0,-0.9) {B: ``Mallory is a CA''};
      \node[digest, font=\small] (d) at (4.2,0) {the same MD5};
      \draw[data] (a.east) -- (d.west); \draw[data] (b.east) -- (d.west);
      \node[asig, font=\small] (s) at (8.0,0.9) {\faStamp\ signature on A};
      \draw[data] (d.east) -- (s.west);
      \node[font=\small, anchor=west, align=left] at (6.9,-0.7) {fits B too \ok\\[2pt]\scriptsize\textcolor{gray}{verify only sees the digest}};
    \end{tikzpicture}
  \end{center}
  \vfill
  \begin{center}\small
    Mallory makes a harmless A and a harmful B with the same digest, gets A signed, then uses the signature with B.\\[4pt]
    Real: a rogue certificate authority (2008), the Flame malware (2012), both with MD5 (Lecture 7).
  \end{center}
  \keymsg{A signature is only as strong as its hash: sign with SHA-256, never MD5 or SHA-1.}
  \say{Podpis obejmuje tylko skrót, więc kolizja skrótu to fałszerstwo podpisu. Oba przypadki z wykładu 7: 2008 badacze zrobili fałszywy urząd certyfikacji (CA, pojęcie z wykładu 9; tu wystarczy: ktoś, kto podpisuje cudze klucze), 2012 Flame miał fałszywy podpis kodu Microsoftu. Dlatego przeglądarki od lat odrzucają podpisy z MD5 i SHA-1. Zadanie dokładnie o tym, na prawdziwym certyfikacie, jest w wykładzie 12 (Task 5); tu go nie powtarzamy.}
\end{frame}

\begin{frame}{See it: RSA signatures}
  \begin{center}\small Textbook RSA on the key \{7, 187\} and \{23, 187\}, the forgeries, and the fix with a hash.\end{center}
  \vfill
  \opensim{signatures/\#2}
  \vfill
  \begin{enumerate}\small
    \item Step 2: change $M$ and watch $S$ and $S^{e}$; all 187 messages come back.
    \item Step 3: pick two signed messages and multiply; pick any $S$ and see the message it signs.
    \item Step 4: the same tricks after SHA-256: the number they sign belongs to no message.
  \end{enumerate}
  \vfill
  \keymsg{The forgeries work on numbers; the hash takes the numbers away from Mallory.}
  \say{Pokaz na żywo, około 3 minut, można pominąć przy braku czasu. Liczby w symulatorze są te same co na slajdach: 11 -> 88, 2 i 3 -> 162 i 181 -> 150, S = 100 -> M = 144.}
\end{frame}

% ======================================================================
\section[What signatures give]{What signatures give, and where}
% ======================================================================

\begin{frame}{Three guarantees}
  \vfill
  \begin{center}
    \begin{tikzpicture}
      \foreach \x/\ic/\name/\txt [count=\k] in {
          0/\faEquals/integrity/{the message was not\\changed after signing},
          4.6/\faUserCheck/authenticity/{it was signed with\\this private key},
          9.2/\faGavel/non-repudiation/{the signer cannot later\\deny it to anybody}} {
        \node[hw, minimum width=4.2cm, minimum height=3.2cm] (b\k) at (\x,0) {};
        \node[text=nInk, font=\fontsize{28}{30}\selectfont] at (\x,0.7) {\ic};
        \node[font=\normalsize\bfseries, text=nInk] at (\x,-0.25) {\name};
        \node[font=\small, text=nInk, align=center] at (\x,-0.95) {\txt};
      }
    \end{tikzpicture}
  \end{center}
  \vfill
  \begin{center}\small A MAC gives the first two, but only to the holders of the key. The third one only a signature gives.\end{center}
  \keymsg{A signature gives integrity, authenticity and non-repudiation.}
  \say{Trzy słowa, które pojawiają się w każdym podręczniku i normie. Integralność: dane się nie zmieniły (punkt 1 na slajdzie What True means). Autentyczność: kto podpisał, a dokładniej: którym kluczem (punkt 2). Niezaprzeczalność (non-repudiation): następny slajd.}
\end{frame}

\samedot
\begin{frame}{Non-repudiation}
  \vfill
  \begin{center}
    \begin{tikzpicture}[>=Stealth]
      \alice{0,0}
      \node[msg, font=\small, text width=3.6cm, align=center] (n) at (4.2,0.6) {``Alice owes Bob 500 pounds.''};
      \node[asig, font=\small] (s) at (4.2,-0.5) {\faStamp\ Alice's signature};
      \carol{10.0,0}
      \node[apub, font=\small] (pk) at (10.0,1.6) {\faKey\ Alice's public key};
      \draw[data] (n.east) -- (C.west);
      \draw[data] (s.east) -- (C.west);
      \draw[data] (pk) -- (C);
      \node[font=\small, anchor=north] at (10.0,-1.3) {\ok\ only Alice's key};
      \node[font=\small, text=sPinkText, anchor=north, align=center] at (0,-1.3) {``I never\\signed that!''};
    \end{tikzpicture}
  \end{center}
  \vfill
  \begin{center}\small
    Carol checks the signature herself. Bob could not have made it: he does not have Alice's private key.\\[2pt]
    \textcolor{gray}{Unless the private key was stolen: the owner must protect it and report a leak.}
  \end{center}
  \keymsg{Non-repudiation: a third party can check, so the signer cannot deny.}
  \say{Porównać ze slajdem Who could have made the tag: z MAC Carol nie wiedziała, czy notatkę zrobiła Alice, czy Bob. Z podpisem wie: tylko klucz prywatny Alice daje podpis, który przechodzi z jej kluczem publicznym. Zastrzeżenie: jeśli klucz ukradziono, Alice może się bronić. Dlatego prawo (w UE rozporządzenie eIDAS) wymaga do podpisu kwalifikowanego klucza w bezpiecznym urządzeniu (karta, token).}
\end{frame}

\samedot
\begin{frame}{Not given: secrecy}
  \vfill
  \begin{center}
    \begin{tikzpicture}[>=Stealth]
      \alice{0,0}
      \node[msg, font=\small] (m) at (3.4,0.45) {``Pay Bob 10 pounds.''};
      \node[asig, font=\small] (s) at (3.4,-0.45) {\faStamp\ signature};
      \bob{10.4,0}
      \draw[net] (5.2,0) -- (B.west);
      \draw[data] (A.east) -- (1.9,0);
      \eve{7.4,-1.6}
      \draw[data, dashed] (7.4,0) -- (E.north);
      \node[font=\small, text=sPinkText, anchor=west] at (8.1,-1.6) {reads everything};
    \end{tikzpicture}
  \end{center}
  \vfill
  \begin{center}\small
    The message travels in the clear next to its signature: Eve reads it. A signature protects against changes, not against reading.\\[4pt]
    For both: sign, then encrypt (or use a protocol that does both, as the web does).
  \end{center}
  \keymsg{A signature is not encryption: the signed message stays readable.}
  \say{Częsty błąd studentów: ,,podpisane, więc tajne''. Nie: podpis jest obok wiadomości, a wiadomość jest jawna (w podręcznikowym RSA z zadania 3 wręcz wyszła z podpisu). Jeśli trzeba obu, łączymy: podpis daje autentyczność, szyfrowanie tajność. W przeglądarce robi to protokół: podpis serwera przy uzgadnianiu klucza, potem szyfrowanie z tagiem.}
\end{frame}

\samedot
\begin{frame}{Not given: freshness}
  \vfill
  \begin{center}
    \begin{tikzpicture}[>=Stealth]
      \foreach \y/\who/\nm in {0.7/Monday/a, -0.7/Tuesday/b} {
        \node[msg, font=\small] (m\nm) at (0,\y) {``Pay Bob 10 pounds.''};
        \node[asig, font=\small] at (2.2,\y) {\faStamp};
        \node[font=\small, text=nInk, anchor=west] at (2.8,\y) {\who: \ok\ valid};
      }
      \mallory{-3.3,-0.7}
      \draw[data] (M.east) -- (mb.west);
      \node[font=\scriptsize, text=sPinkText] at (-3.3,0.5) {sends it again};
    \end{tikzpicture}
  \end{center}
  \vfill
  \begin{center}\small
    A copied signature stays valid forever: Mallory can send the same signed order again (a \textbf{replay}).\\[4pt]
    The fix is in the message: a date, a number used once, or a counter, signed together with the text.
  \end{center}
  \keymsg{A signature says what was signed, not when or how many times.}
  \say{Atak powtórzenia (replay): Mallory nie zmienia ani bajtu, więc podpis przechodzi. Ochrona: w podpisanej treści musi być coś jednorazowego (numer zlecenia, data, liczba losowa od odbiorcy). Bank odrzuca zlecenie, którego numer już widział. W wykładzie 11 zobaczymy, że serwer podpisuje cały przebieg połączenia z liczbami losowymi obu stron, więc stary podpis nie pasuje do nowego połączenia.}
\end{frame}

\begin{frame}{MAC or signature}
  \vfill
  \begin{center}
    \renewcommand{\arraystretch}{1.5}
    \begin{tabular}{@{}l l l@{}}
      \toprule
      & \textbf{MAC} (HMAC-SHA256) & \textbf{signature} (Ed25519) \\
      \midrule
      key & one \keyw{shared} key & a key pair \\
      who can make it & everybody with the key & only the private key owner \\
      who can check it & only those with the key & everybody (public key) \\
      convinces a third party & \nok\ no & \ok\ yes \\
      size & 32 bytes & 64 bytes \\
      speed & very fast & slower \\
      \bottomrule
    \end{tabular}
  \end{center}
  \vfill
  \keymsg{A MAC is for two parties who share a key; a signature is for everybody else.}
  % CHECKED: HMAC-SHA256 tag 32 bytes (RFC 2104 with SHA-256); Ed25519 signature 64 bytes (RFC 8032); speed: Task 7 and hmac in python3
  \say{Porównanie. Szybkość: HMAC to dwa skróty SHA-256, kilka mikrosekund dla krótkiej wiadomości; podpis Ed25519 około 0,1 ms, RSA-2048 ponad 1 ms (pomiar w zadaniu 7). Najważniejszy wiersz: przekonuje osobę trzecią.}
\end{frame}

\samedot
\begin{frame}{When a MAC is the better choice}
  \vfill
  \begin{columns}[t,onlytextwidth]
    \column{0.32\textwidth}
      \stepbox{\faHandshake\ Both ends share a key anyway}{After a key exchange (Lecture 4), each data packet of a web connection carries a tag, not a signature.}
    \column{0.32\textwidth}
      \stepbox{\faServer\ One party checks its own data}{A server tags the cookie it gives you and checks it when it comes back.}
    \column{0.32\textwidth}
      \stepbox{\faUserSecret\ Deniability is wanted}{Off-the-Record messaging (2004) uses MACs on purpose: the chat cannot later be proved to a third party.}
  \end{columns}
  \vfill
  \keymsg{``Cannot convince a third party'' is sometimes exactly what we want.}
  % CHECKED: Borisov, Goldberg, Brewer, "Off-the-Record Communication, or, Why Not To Use PGP", WPES 2004 (cs.uwaterloo.ca/~iang/pubs/otr-wpes.pdf, 2026-11-17): messages authenticated with MACs, MAC keys published later
  \say{Ważne, żeby nie przedstawić MAC jako gorszego. W połączeniu przeglądarki dane po uzgodnieniu klucza są chronione tagiem (np. ChaCha20-Poly1305 z wykładu 2): podpis każdego pakietu byłby wolny i niepotrzebny, bo obie strony mają klucz. Ciasteczko: serwer sprawdza swoje własne dane, nikt trzeci nie musi. Deniability: komunikator Off-the-Record (Borisov, Goldberg, Brewer, 2004) celowo używa MAC, a klucze MAC po rozmowie publikuje, żeby nikt nie mógł potem udowodnić, kto co napisał; to cecha, nie wada.}
\end{frame}

\begin{frame}{Task 6: MAC or signature?}
  \vfill
  \small
  For each situation, choose a MAC or a signature, and say why in one line.
  \begin{enumerate}\setlength{\itemsep}{3pt}
    \item A phone app checks that an update really comes from its maker.
    \item Two servers of one company pass messages to each other many times per second.
    \item A student hands in a thesis; years later, the university must prove which file was handed in.
    \item A web server gives you a ``logged in'' cookie and checks it on each visit.
    \item Alice and Bob chat and do not want their chat to be provable to anyone else.
  \end{enumerate}
  \vfill
  \onpaper{about 6 minutes}
  \vfill
  \say{Na kartce, około 6 minut, można w parach. Odpowiedzi na następnym slajdzie. Jeśli brakuje czasu: omówić wspólnie, każdy punkt po 30 sekund. Kluczowe pytania pomocnicze: ile stron sprawdza? Czy trzeba przekonać kogoś trzeciego? Czy strony i tak mają wspólny klucz?}
\end{frame}

\samedot
\begin{frame}{Task 6: what we saw}
  \vfill
  \begin{center}
    \renewcommand{\arraystretch}{1.45}
    \begin{tabular}{@{}c l l@{}}
      \toprule
      & \textbf{choice} & \textbf{why} \\
      \midrule
      1 & signature & millions check; none of them may be able to make updates \\
      2 & MAC & two parties, one owner, a shared key is easy, speed matters \\
      3 & signature & a third party (a court, an auditor) must be convinced, years later \\
      4 & MAC & the server checks its own tag; nobody else checks \\
      5 & MAC & deniability: a tag does not prove who wrote it \\
      \bottomrule
    \end{tabular}
  \end{center}
  \vfill
  \keymsg{Many checkers or a third party: a signature. Two parties with a shared key: a MAC.}
  \say{Odpowiedzi. W punkcie 2 podpis też zadziała, ale bez potrzeby jest wolniejszy; w punkcie 3 dodatkowo przyda się znacznik czasu. Jeśli ktoś wybrał inaczej i dobrze uzasadnił, to też jest dobra odpowiedź: liczy się uzasadnienie.}
\end{frame}

\begin{frame}{Signatures on curves}
  \vfill
  \begin{center}
    \begin{tikzpicture}
      \foreach \x/\name/\curve/\used [count=\k] in {0/ECDSA/{the curve P-256}/{TLS, Bitcoin, passports,\\the PlayStation 3}, 7/Ed25519/{Curve25519 (Lecture 5)}/{SSH, Signal, Git,\\many new systems}} {
        \node[hw, minimum width=6.2cm, minimum height=3.0cm] (b\k) at (\x,0) {};
        \node[font=\large\bfseries, text=nInk] at (\x,0.9) {\name};
        \node[font=\small, text=nInk] at (\x,0.3) {\curve};
        \node[font=\small, text=gray, align=center] at (\x,-0.6) {\used};
      }
    \end{tikzpicture}
  \end{center}
  \vfill
  \begin{center}\small
    The same sign and verify, keys of 32 bytes instead of 256: the shorter keys of Lecture 5.\\
    We keep both as black boxes, with one exception: the random number of ECDSA.
  \end{center}
  \keymsg{On curves: short keys, short signatures, the same interface.}
  % CHECKED: FIPS 186-5 (2023) specifies RSA, ECDSA and EdDSA (csrc.nist.gov/Projects/digital-signatures, 2026-11-17)
  \say{ECDSA to wersja podpisu ElGamala/DSA na krzywej; standard NIST (FIPS 186, od 2000 roku; FIPS 186-5 z 2023 dodał też EdDSA, czyli Ed25519). Przykłady użycia: certyfikaty i TLS (bardzo często ECDSA P-256), Bitcoin (ECDSA na krzywej secp256k1), paszporty biometryczne, konsola PS3 (za chwilę). Ed25519: klucze SSH, Signal, podpisy commitów. Matematyki w środku nie omawiamy, poza jednym wzorem potrzebnym do ataku.}
\end{frame}

\samedot
\begin{frame}{ECDSA needs a fresh random k}
  \vfill
  \begin{center}
    \begin{tikzpicture}[>=Stealth]
      \node[msg, font=\small] (m) at (0,0.9) {\faFile*[regular]\ message};
      \node[apriv, font=\small] (sk) at (0,0) {\faKey\ private key $d$};
      \node[hw, font=\small, draw=AccentOrange, line width=1.2pt] (k) at (0,-0.9) {\faDice\ random $k$, new};
      \node[hwmain, minimum width=2.2cm, minimum height=2.2cm] (s) at (3.6,0) {\Large\faCogs\\ \small\bfseries ECDSA sign};
      \draw[data] (m.east) -- (s.west |- m.east);
      \draw[data] (sk.east) -- (s.west |- sk.east);
      \draw[data] (k.east) -- (s.west |- k.east);
      \node[asig, font=\small] (sig) at (7.2,0) {\faStamp\ $(r, s)$};
      \draw[data] (s) -- (sig);
    \end{tikzpicture}
  \end{center}
  \vfill
  \begin{itemize}\small
    \item Every ECDSA signature needs a new random number $k$, like ElGamal in Lecture 4.
    \item The same message signed twice gives two different signatures; both are valid.
    \item $k$ must stay \textbf{secret} and must \textbf{never repeat}. We will see why.
  \end{itemize}
  \keymsg{ECDSA: a fresh secret random number for every signature.}
  \say{Podpis ECDSA to dwie liczby, r i s (po 32 bajty dla P-256). k jest jak k w ElGamalu z wykładu 4: losowe, tajne, nowe dla każdej wiadomości. To jest najsłabszy punkt ECDSA w praktyce: dobra losowość w małym urządzeniu nie jest oczywista.}
\end{frame}

\samedot
\begin{frame}{Ed25519 makes its own k}
  \vfill
  \begin{center}
    \begin{tikzpicture}[>=Stealth]
      \node[apriv, font=\small] (sk) at (0,0.6) {\faKey\ private key};
      \node[msg, font=\small] (m) at (0,-0.6) {\faFile*[regular]\ message};
      \mech{h}{3.4,0}{SHA-512}{2.2cm}
      \node[hw, font=\small] (k) at (6.6,0) {$k$};
      \draw[data] (sk.east) -- (h.west); \draw[data] (m.east) -- (h.west); \draw[data] (h) -- (k);
      \node[font=\small, text=nInk, anchor=west, align=left] at (7.4,0) {the same message $\to$ the same $k$\\another message $\to$ another $k$};
    \end{tikzpicture}
  \end{center}
  \vfill
  \begin{center}\small
    Ed25519 computes $k$ from the private key and the message: no random number generator,\\
    so a bad generator cannot hurt it. The same idea for ECDSA: deterministic $k$, RFC 6979 (2013).
  \end{center}
  \keymsg{Deterministic $k$: secret, different for every message, never from a random generator.}
  % CHECKED: RFC 8032 section 5.1.6 (r = SHA-512(prefix || M), prefix from the hash of the private key); RFC 6979 (August 2013)
  \say{Dlatego Ed25519 jest deterministyczny (zadanie 2: ten sam podpis dwa razy). k jest tajne, bo zależy od klucza prywatnego, i różne dla różnych wiadomości, bo zależy od wiadomości. Te same wiadomości dają to samo k, ale wtedy też ten sam podpis, więc nic nie wycieka. RFC 6979 to ta sama sztuczka dla ECDSA.}
\end{frame}

\begin{frame}{Task 7: three schemes}
  \vfill
  \taskgoal{compare RSA-2048, ECDSA P-256 and Ed25519: sizes, the same message signed twice, and speed.}
  \taskready{\texttt{schemes.py}: \texttt{load\_keys}, \texttt{sign}, \texttt{verify}, \texttt{public\_key\_size}; Alice's three keys in PEM files.}
  \vfill
  \opentask{Task7}{about 8 minutes}
  \vfill
  \say{Około 8 minut. Jedna pętla po trzech nazwach. TODO 1: rozmiary. TODO 2: dwa podpisy tej samej wiadomości: równe? oba poprawne? TODO 3: czas 200 podpisów i 200 sprawdzeń (time.perf\_counter). Przed TODO 2 zapytać: który da dwa takie same podpisy? RSA-2048 w module to RSA-PSS (z solą). Przy braku czasu pominąć TODO 3 i pokazać pomiar z następnego slajdu.}
\end{frame}

\samedot
\begin{frame}{Task 7: what we saw}
  % CHECKED: python3 Task7/solution/three_schemes.py, 2026-11-17, Intel Core i5-8265U, OpenSSL 3.0.13, cryptography 41.0.7
  %   (three runs; times vary by about 20%; the first run is shown)
  \vfill
  \begin{center}
    \renewcommand{\arraystretch}{1.5}
    \begin{tabular}{@{}l r r c r r@{}}
      \toprule
      & \textbf{public key} & \textbf{signature} & \textbf{signed twice: same?} & \textbf{sign} & \textbf{verify} \\
      \midrule
      RSA-2048 (PSS) & 256 B & 256 B & no (random salt) & 1.08 ms & 0.09 ms \\
      ECDSA P-256 & 33 B & 64 B & no (random $k$) & 0.07 ms & 0.14 ms \\
      Ed25519 & 32 B & 64 B & \textbf{yes} & 0.08 ms & 0.21 ms \\
      \bottomrule
    \end{tabular}
  \end{center}
  \vfill
  \begin{center}\small Measured on a laptop (Intel Core i5-8265U, OpenSSL 3.0.13). All signatures were valid.\end{center}
  \keymsg{Curves: 8 times smaller keys and signatures, and signing 15 times faster than RSA.}
  \say{Wynik pierwszego uruchomienia; u studentów liczby będą inne, proporcje podobne. Klucz ECDSA 33 bajty: punkt w postaci skróconej (x i jeden bit y). Podpis ECDSA tu to r i s po 32 bajty; w formacie DER (np. w certyfikacie) ma około 70 do 72 bajtów. RSA weryfikuje bardzo szybko (małe e = 65537), ale podpisuje wolno (duże d). Dlatego serwery, które podpisują przy każdym połączeniu, coraz częściej mają klucze ECDSA.}
\end{frame}

\begin{frame}{Recall: k used twice}
  \vfill
  \begin{center}
    \begin{tikzpicture}
      \node[anchor=east, font=\small, text=nInk] at (-0.2,0.5) {first:};
      \node[hw, font=\small, draw=AccentOrange, very thick] at (1.4,0.5) {$C_1 = g^{k}$};
      \node[hw, font=\small] at (4.6,0.5) {$C_2 = M_1 \times y^{k}$};
      \node[anchor=east, font=\small, text=nInk] at (-0.2,-0.5) {second:};
      \node[hw, font=\small, draw=AccentOrange, very thick] at (1.4,-0.5) {$C_1 = g^{k}$};
      \node[hw, font=\small] at (4.6,-0.5) {$C_2 = M_2 \times y^{k}$};
      \node[font=\scriptsize\bfseries, text=AccentOrange] at (1.4,-1.2) {the same $k$, visible to all};
    \end{tikzpicture}
  \end{center}
  \vfill
  \begin{center}\small
    Lecture 4, ElGamal: with the same $k$, the same mask $y^{k}$ hid two messages.\\
    One known message gave the mask, and the mask gave the other message.
  \end{center}
  \keymsg{A repeated secret random number leaks; ECDSA has the same weakness.}
  \say{Przypomnienie z wykładu 4 (zadanie 6, k used twice). Tam z powtórzonego k wyciekała wiadomość. W ECDSA wycieknie coś gorszego: klucz prywatny.}
\end{frame}

\samedot
\begin{frame}{Inside ECDSA: r and s}
  \vfill
  \stepbox{1. A point from k}{$R = k\,G$ \quad and \quad $r$ = the $x$ of $R$, mod $q$ \hfill \textcolor{gray}{(G and its order $q$: Lecture 5)}}
  \stepbox{2. Mix in the digest and the private key}{$s = (z + r\,d) \,/\, k \pmod q$ \hfill \textcolor{gray}{($z$ = the digest as a number, $d$ = the private key)}}
  \stepbox{3. The signature}{$(r, s)$. Verify uses only $Q = d\,G$, never $d$ or $k$.}
  \vfill
  \begin{center}\small
    One equation, two unknowns $k$ and $d$: as long as $k$ is secret and new, $s$ hides $d$.
  \end{center}
  \curious{signatures/\#5}
  \keymsg{$s$ hides the private key $d$ behind the secret $k$.}
  \say{Tylko tyle wzorów, ile trzeba do ataku. G: punkt startowy, q: rząd G (wykład 5, slajd The order of G: po q krokach wracamy do O), wszystkie liczby bierzemy mod q. Dzielenie mod q to mnożenie przez odwrotność, jak w wykładzie 3. Jak verify sprawdza bez d: liczy punkt z G i Q i porównuje jego x z r; szczegóły w symulatorze, krok 5. Dlaczego bezpieczne: dla jednego podpisu jest jedno równanie i dwie niewiadome (k i d), a k nie da się wyciągnąć z R (problem ECDLP z wykładu 5).}
\end{frame}

\samedot
\begin{frame}{The example on our curve}
  % CHECKED: python3 (scratchpad l8/toy.py): curve y^2 = x^3 + 2x + 2 mod 17, G = (5, 1), q = 19; d = 7, Q = 7G = (0, 6);
  %   k = 10, 10G = (7, 11), r = 7; z = 5: s = 13; z = 11: s = 6; both verify
  \vfill
  \begin{columns}[c,onlytextwidth]
    \column{0.46\textwidth}
      \centering
      \modgrid[0.2]{\gmark[nInk]{5,1}{$G$}{below right}\gmark[sIndigo]{0,6}{$Q$}{right}\gmark[AccentOrange]{7,11}{$kG$}{above right}}
    \column{0.52\textwidth}
      \small
      Curve of Lecture 5: $y^2 = x^3 + 2x + 2 \bmod 17$, $G = (5, 1)$, $q = 19$.\\[6pt]
      Alice: $d = 7$, \;$Q = 7G = (0, 6)$.\\[6pt]
      \textbf{Signature 1}, $z_1 = 5$, \;$k = 10$:\\
      \quad $10G = (7, 11)$, so $r = 7$\\
      \quad $s_1 = (5 + 7 \times 7) / 10 = \mathbf{13} \pmod{19}$\\[6pt]
      \textbf{Signature 2}, $z_2 = 11$, the \textbf{same} $k = 10$:\\
      \quad $r = 7$ again, \;$s_2 = (11 + 49) / 10 = \mathbf{6} \pmod{19}$
  \end{columns}
  \vfill
  \keymsg{The same $k$ gives the same $r$: anybody can see that $k$ repeated.}
  \say{Malutki przykład na krzywej z wykładu 5 (19 punktów razem z O). Hashe z1 = 5 i z2 = 11 to zabawkowe liczby mod 19 zamiast prawdziwego SHA-256. Rachunek: 5 + 49 = 54 = 16 mod 19; odwrotność 10 mod 19 to 2 (10 razy 2 = 20 = 1), więc s1 = 16 razy 2 = 32 = 13. Drugi: 11 + 49 = 60 = 3 mod 19, razy 2 = 6. Oba podpisy przechodzą sprawdzenie (policzone w Pythonie). Najważniejsze: r jest w obu takie samo, więc każdy widzi, że k się powtórzyło.}
\end{frame}

\samedot
\begin{frame}{The same k twice, step 1: find k}
  \vfill
  \stepbox{Two signatures, the same $k$ and $r$}{$s_1 = (z_1 + r\,d) / k$ \qquad $s_2 = (z_2 + r\,d) / k$}
  \stepbox{Subtract: $r\,d$ disappears}{$s_1 - s_2 = (z_1 - z_2) / k$ \quad $\Rightarrow$ \quad $k = (z_1 - z_2) / (s_1 - s_2) \pmod q$}
  \vfill
  \begin{center}
    Our curve: \quad $k = (5 - 11) / (13 - 6) = -6 / 7 = 13 \times 11 = 143 = \mathbf{10} \pmod{19}$
  \end{center}
  \vfill
  \keymsg{Two equations, and the private key cancels out: k is no longer secret.}
  % CHECKED: python3: (-6) % 19 = 13, pow(7, -1, 19) = 11, 13 * 11 % 19 = 10
  \say{z1 i z2 to skróty dwóch podpisanych wiadomości: każdy je policzy, bo wiadomości są znane. s1, s2 i r są w podpisach. Odejmujemy równania i r razy d znika, więc zostaje jedna niewiadoma: k. Rachunek: -6 mod 19 = 13; odwrotność 7 mod 19 to 11 (7 razy 11 = 77 = 4 razy 19 + 1); 13 razy 11 = 143 = 7 razy 19 + 10, czyli k = 10, zgadza się.}
\end{frame}

\samedot
\begin{frame}{The same k twice, step 2: find d}
  \vfill
  \stepbox{From the first signature}{$s_1 = (z_1 + r\,d) / k$ \quad $\Rightarrow$ \quad $d = (s_1\,k - z_1) / r \pmod q$}
  \vfill
  \begin{center}
    Our curve: \quad $d = (13 \times 10 - 5) / 7 = 125 / 7 = 11 \times 11 = 121 = \mathbf{7} \pmod{19}$ \;\ok
  \end{center}
  \vfill
  \begin{center}\small Mallory now has Alice's private key: she can sign anything as Alice.\end{center}
  \keymsg{A repeated $k$ in ECDSA gives away the private key, from two public signatures.}
  % CHECKED: python3: 125 % 19 = 11, pow(7, -1, 19) = 11, 11 * 11 % 19 = 7
  \say{Jedno równanie, jedna niewiadoma d. 125 mod 19 = 11 (125 = 6 razy 19 + 11), razy odwrotność 7, czyli 11: 121 mod 19 = 7. Zgadza się z kluczem Alice. To jest gorsze niż w ElGamalu: tam wyciekała jedna wiadomość, tu cały klucz prywatny, czyli wszystkie przyszłe podpisy. Symulator, krok 5, liczy ten atak krok po kroku (pokażemy po zadaniu 8).}
\end{frame}

\begin{frame}{Sony PlayStation 3, 2010}
  \vfill
  \begin{columns}[c,onlytextwidth]
    \column{0.3\textwidth}
      \centering
      \begin{tikzpicture}
        \node[text=nInk, font=\fontsize{60}{60}\selectfont] {\faGamepad};
      \end{tikzpicture}
    \column{0.68\textwidth}
      \stepbox{The design}{The console runs only software signed with Sony's ECDSA private key.}
      \stepbox{The mistake}{Sony's signing used the \textbf{same $k$} for every signature.}
      \probbox{The result}{In December 2010 the group fail0verflow showed at the 27th Chaos Communication Congress that two signatures give the private key. Anybody could sign software as Sony.}
  \end{columns}
  \vfill
  \keymsg{A strong curve and a correct formula, broken by one number that did not change.}
  % CHECKED: fail0verflow, "Console Hacking 2010", 27C3, Berlin, 29 December 2010 (exophase.com/20540, siliconangle.com/?p=35925, games.slashdot.org 2010-12-29; 2026-11-17)
  \say{Prawdziwa historia: na kongresie 27C3 w Berlinie (grudzień 2010) grupa fail0verflow pokazała, że Sony w podpisach programów dla PS3 używało stałego k. Z dwóch podpisanych plików wyliczyli klucz prywatny; kilka dni później klucz został opublikowany w internecie. Kryptografia była poprawna, implementacja nie. Sony nie mogło zmienić klucza w milionach konsol, które już były u ludzi.}
\end{frame}

\begin{frame}{Task 8: the same k twice}
  \vfill
  \taskgoal{from two console programs signed with the same $k$, recover the private key, then sign Mallory's program.}
  \taskready{\texttt{ecdsa\_ready.py}: \texttt{q}, \texttt{hash\_number}, \texttt{public\_key}, \texttt{sign}, \texttt{verify} on the real curve P-256.}
  \vfill
  \opentask{Task8}{about 12 minutes}
  \vfill
  \say{Około 12 minut. Prawdziwa krzywa P-256 i SHA-256, liczby po 77 cyfr. Najważniejsze są TODO 2 i 3: dwa wzory ze slajdów step 1 i step 2, z dzieleniem jako pow(x, -1, q). TODO 4 sprawdza, czy public\_key(d) == Q. TODO 5: podpis programu Mallory odzyskanym kluczem, konsola mówi True. Typowy błąd: zapomnieć o \% q albo podzielić zwykłym /. Symulatora z atakiem nie otwierać przed zadaniem.}
\end{frame}

\samedot
\begin{frame}{Task 8: what we saw}
  % CHECKED: python3 Task8/solution/recover_key.py, 2026-11-17
  \vfill
  \begin{center}
    \renewcommand{\arraystretch}{1.45}
    \begin{tabular}{@{}l l@{}}
      \toprule
      $r_1 = r_2$ & \texttt{a7ef236729f03041\ldots 6e3800979}: the same $k$ \\
      $k$ & \texttt{ccfc9b5a4810d392\ldots 2e518abbc0} \\
      $d$ & \texttt{db8c7cb153011659\ldots 2aa1189f1c} \\
      \texttt{public\_key(d) == Q} & \ok\ \texttt{True} \\
      the console accepts Mallory's program & \ok\ \texttt{True} \\
      \bottomrule
    \end{tabular}
  \end{center}
  \vfill
  \begin{center}\small Two formulas and two inverses mod $q$: the private key of the console maker.\end{center}
  \keymsg{Mallory signs as Alice Games, and every console believes her.}
  \say{Wynik rozwiązania. k = (z1 - z2) razy odwrotność (s1 - s2) mod q, d = (s1 razy k - z1) razy odwrotność r mod q. Całość działa w ułamku sekundy. Konsola sprawdza podpis programu Mallory kluczem publicznym Alice Games i mówi True, bo to jest prawdziwy podpis zrobiony prawdziwym kluczem prywatnym.}
\end{frame}

\begin{frame}{The rule for k}
  \vfill
  \begin{columns}[t,onlytextwidth]
    \column{0.48\textwidth}
      \centering
      \begin{tikzpicture}
        \node[text=nInk, font=\fontsize{44}{44}\selectfont] {\faDice};
      \end{tikzpicture}\\[6pt]
      \textbf{Random $k$}: secret, new for every signature,\\ from a good random generator.
    \column{0.48\textwidth}
      \centering
      \begin{tikzpicture}
        \node[text=nInk, font=\fontsize{44}{44}\selectfont] {\faHashtag};
      \end{tikzpicture}\\[6pt]
      \textbf{Deterministic $k$}: computed from the private key and the message (Ed25519, RFC 6979).
  \end{columns}
  \vfill
  \begin{center}\small Not only a repeated $k$: even a few known or biased bits of $k$ over many signatures can leak $d$.\end{center}
  \keymsg{The same lesson as the keystream, the nonce and ElGamal: a value meant to be used once must be used once.}
  \say{Gdyby ktoś pytał: nawet częściowo przewidywalne k (kilka bitów znanych w wielu podpisach) pozwala odzyskać klucz metodami kratowymi (lattice); to był problem m.in. kart i bibliotek z błędnym generatorem. Dlatego dziś zaleca się deterministyczne k albo Ed25519. Ta sama lekcja wróciła w kursie po raz czwarty: keystream (wykład 1), nonce (wykład 2), k w ElGamalu (wykład 4), k w ECDSA.}
\end{frame}

\samedot
\begin{frame}{Good to know: Bitcoin wallets on Android, 2013}
  \vfill
  \begin{columns}[c,onlytextwidth]
    \column{0.3\textwidth}
      \centering
      \begin{tikzpicture}
        \node[text=nInk, font=\fontsize{56}{56}\selectfont] {\faBitcoin};
      \end{tikzpicture}
    \column{0.68\textwidth}
      \stepbox{The bug}{A flaw in Android's random number generator: some wallet apps signed two payments with the \textbf{same $k$}.}
      \probbox{The result}{All payments are public in the blockchain. Attackers searched it for two signatures with the same $r$, computed the private keys and took the coins (bitcoin.org alert, 11 August 2013).}
  \end{columns}
  \vfill
  \keymsg{The PlayStation attack, three years later, with real money.}
  % CHECKED: bitcoin.org/en/alert/2013-08-11-android (via mailing list of M. Hearn), The Register 2013-08-12, Symantec blog 2013-08 (2026-11-17)
  \say{Można pominąć. 11 sierpnia 2013 bitcoin.org ogłosił, że portfele bitcoin na Androidzie są zagrożone: systemowy generator SecureRandom miał błąd i czasem dawał to samo k. Ponieważ wszystkie transakcje są publiczne, wystarczyło przeszukać łańcuch bloków w poszukiwaniu dwóch podpisów z tym samym r i policzyć klucz, jak w zadaniu 8. Poprawione portfele czytały losowość bezpośrednio z /dev/urandom.}
\end{frame}

\begin{frame}{See it: ECDSA with a repeated k}
  \begin{center}\small ECDSA on the small curve of Lecture 5, with $q = 19$: sign, verify, and the attack.\end{center}
  \vfill
  \opensim{signatures/\#5}
  \vfill
  \begin{enumerate}\small
    \item Step 5.1: choose $d$, $k$ and $z$; watch $R = kG$, $r$ and $s$; verify draws the point it checks.
    \item Step 5.2: sign two messages with the same $k$; the two $r$ are equal.
    \item Press ``recover'': $k$ and then $d$ appear, step by step, as on the slides.
  \end{enumerate}
  \vfill
  \keymsg{The same $r$ twice is a public alarm: the private key can be computed.}
  \say{Pokaz na żywo, około 3 minut. Domyślne liczby to przykład ze slajdów: d = 7, k = 10, z = 5 i 11, podpisy (7, 13) i (7, 6).}
\end{frame}

\begin{frame}{Where signatures are used}
  \vfill
  \begin{center}
    \begin{tikzpicture}
      \foreach \ic/\name/\sub [count=\k from 0] in {
          \faDownload/software updates/{the update is checked\\before it runs},
          \faCodeBranch/signed commits/{who wrote this change\\in the repository},
          \faFileSignature/documents/{contracts, invoices,\\a thesis as a PDF},
          \faLock/the web/{the server proves\\who it is}} {
        \pgfmathsetmacro{\x}{6.9*mod(\k,2)}
        \pgfmathsetmacro{\y}{-2.3*int(\k/2)}
        \node[hw, minimum width=6.6cm, minimum height=2.0cm] (b) at (\x,\y) {};
        \node[text=nInk, font=\fontsize{24}{26}\selectfont, anchor=west] at ([xshift=0.25cm]b.west) {\ic};
        \node[font=\small\bfseries, text=nInk, anchor=north west] at ([xshift=1.5cm,yshift=-0.3cm]b.north west) {\name};
        \node[font=\scriptsize, text=nInk, anchor=north west, align=left, inner sep=0pt] at ([xshift=1.5cm,yshift=-0.8cm]b.north west) {\sub};
      }
    \end{tikzpicture}
  \end{center}
  \vfill
  \keymsg{Wherever many must check what one has written.}
  \say{Cztery miejsca, po jednym slajdzie każde. Wspólne: jeden podpisuje, wielu sprawdza, często po latach. Przy braku czasu wystarczy ten slajd i dwa następne.}
\end{frame}

\samedot
\begin{frame}{Signed software updates}
  \vfill
  \begin{center}
    \begin{tikzpicture}[>=Stealth]
      \alice{0,0}
      \node[msg, font=\small] (u) at (3.2,0.5) {\faFile*[regular]\ update};
      \node[asig, font=\small] (s) at (3.2,-0.5) {\faStamp\ signature};
      \draw[data] (A.east) -- (u.west); \draw[data] (A.east) -- (s.west);
      \node[hw, minimum width=3.6cm, minimum height=2.2cm] (dev) at (8.6,0) {};
      \node[font=\small\bfseries, text=nInk, anchor=north] at (dev.north) {\faLaptop\ Bob's computer};
      \node[apub, font=\scriptsize] at (8.6,0.0) {\faKey\ Alice's key, built in};
      \node[font=\scriptsize, text=nInk] at (8.6,-0.7) {runs it only if \texttt{True}};
      \draw[net] (4.6,0) -- (dev.west);
    \end{tikzpicture}
  \end{center}
  \vfill
  \begin{itemize}\small
    \item Linux package managers (e.g.\ \texttt{apt}) check the signature on the list of packages and their digests.
    \item Android installs an app only if it is signed; an update must be signed with the same key.
    \item Windows checks the signature of programs and drivers.
  \end{itemize}
  \keymsg{The public key is built in at installation, so a changed server cannot change it.}
  % CHECKED: Debian SecureApt: InRelease / Release.gpg signed with OpenPGP, Release lists SHA-256 of the package lists (wiki.debian.org/SecureApt);
  %   Android: "Android requires that all APKs be digitally signed ... before they are installed or updated" (source.android.com/docs/security/features/apksigning)
  \say{To jest dokładnie problem z początku wykładu, rozwiązany. Debian i Ubuntu: plik InRelease (lista pakietów z ich skrótami SHA-256) jest podpisany, a apt ma klucze publiczne dystrybucji zainstalowane razem z systemem. Pakiety sprawdza się skrótami z podpisanej listy, więc jeden podpis chroni tysiące plików. Android: każda aplikacja musi być podpisana, a aktualizacja tym samym kluczem co wersja zainstalowana. Windows: Authenticode, podpisy programów i sterowników.}
\end{frame}

\samedot
\begin{frame}{Signed commits}
  \vfill
  \begin{center}
    \begin{tikzpicture}
      \node[hw, text width=9.2cm, align=left, font=\small\ttfamily, inner sep=8pt] (c) {
        commit 3f9a2c1\ldots\\
        Author: Alice <alice@example.org>\\[4pt]
        \ \ \ \ Fix the login check};
      \node[draw=aExec, thick, rounded corners=3pt, font=\scriptsize\bfseries, text=aExec, anchor=north east] at ([xshift=-6pt,yshift=-6pt]c.north east) {\faCheck\ Verified};
    \end{tikzpicture}
  \end{center}
  \vfill
  \begin{itemize}\small
    \item The ``Author'' line is plain text: anybody can write ``Alice'' there.
    \item \texttt{git commit -S} adds a signature with Alice's OpenPGP key or, since Git 2.34 (2021), her SSH key (often Ed25519).
    \item GitHub shows ``Verified'' when the signature fits a public key Alice uploaded to her account.
  \end{itemize}
  \keymsg{A signed commit proves which key approved the change, not just which name was typed.}
  % CHECKED: Git 2.34 (November 2021) adds SSH signing, gpg.format = ssh (github.blog/2021-11-15-highlights-from-git-2-34, 2026-11-17)
  \say{W Git autora wpisuje się w konfiguracji (git config user.name), więc każdy może podpisać się w commicie jako ktokolwiek. Podpis commitu obejmuje cały commit (a przez skróty także całą historię wcześniej). Od Git 2.34 (listopad 2021) można podpisywać kluczem SSH. Weryfikacja na GitHubie znowu opiera się na pytaniu ,,czyj to klucz'': GitHub ufa kluczowi, który Alice wgrała na swoje konto.}
\end{frame}

\samedot
\begin{frame}{Signed documents}
  \vfill
  \begin{columns}[c,onlytextwidth]
    \column{0.3\textwidth}
      \centering
      \begin{tikzpicture}
        \node[text=nInk, font=\fontsize{56}{56}\selectfont] {\faFilePdf[regular]};
        \node[asig, font=\scriptsize] at (0.6,-0.9) {\faStamp};
      \end{tikzpicture}
    \column{0.68\textwidth}
      \stepbox{In the file}{A PDF can carry signatures inside it: the reader shows who signed and whether anything changed since.}
      \stepbox{In law}{In the EU (eIDAS regulation, 2014), a \textbf{qualified} electronic signature has the same legal effect as a handwritten one.}
  \end{columns}
  \vfill
  \keymsg{Non-repudiation in practice: a signature that a court accepts.}
  % CHECKED: Regulation (EU) No 910/2014 (eIDAS), Article 25(2): "A qualified electronic signature shall have the equivalent legal effect of a handwritten signature"
  \say{Podpis kwalifikowany: klucz prywatny w certyfikowanym urządzeniu (karta, zdalny moduł), a klucz publiczny poświadczony przez zaufanego dostawcę. W Polsce to np. podpis kwalifikowany na karcie albo w chmurze; profil zaufany to inny mechanizm. Szczegóły prawne pomijamy, ważne jest, że niezaprzeczalność z tego wykładu działa w sądzie.}
\end{frame}

\samedot
\begin{frame}{Recall: the man in the middle}
  \vfill
  \begin{center}
    \begin{tikzpicture}[>=Stealth]
      \alice{0,0}
      \mallory{5,0}
      \bob{10,0}
      \draw[data] (A.east) to node[above, font=\scriptsize] {$A$} (M.west);
      \draw[data] (M.east) to node[above, font=\scriptsize] {$g^{m}$ instead of $A$} (B.west);
      \draw[data] ([yshift=-8pt]B.west) to node[below, font=\scriptsize] {$B$} ([yshift=-8pt]M.east);
      \draw[data] ([yshift=-8pt]M.west) to node[below, font=\scriptsize] {$g^{m}$ instead of $B$} ([yshift=-8pt]A.east);
    \end{tikzpicture}
  \end{center}
  \vfill
  \begin{center}\small
    Lecture 4: Mallory replaced both Diffie-Hellman values and shared one key with Alice and another with Bob.\\
    The fix promised there: Bob seals $B$ with his long-term private key.
  \end{center}
  \keymsg{Diffie-Hellman alone cannot tell who sent a value; a signature can.}
  \say{Przypomnienie ataku z wykładu 4 (zadanie z PIN-em). Pieczęć z tamtego wykładu to dzisiejszy podpis: Bob podpisuje B swoim długoterminowym kluczem prywatnym, a Alice sprawdza go jego kluczem publicznym. Mallory nie umie podpisać g do m jako Bob.}
\end{frame}

\samedot
\begin{frame}{Signatures on the web}
  \vfill
  \begin{center}
    \begin{tikzpicture}[>=Stealth]
      \node[hw, font=\small, minimum width=3.0cm, minimum height=1.2cm] (b) at (0,0) {\faGlobe\ browser};
      \node[hw, font=\small, minimum width=3.0cm, minimum height=1.2cm] (s) at (9,0) {\faServer\ server};
      \draw[data, <->] (b) to node[above, font=\scriptsize] {key exchange (Lecture 4)} (s);
      \node[asig, font=\scriptsize, text width=4.2cm, align=center] at (4.5,-1.1) {\faStamp\ the server signs the exchange with its private key};
    \end{tikzpicture}
  \end{center}
  \vfill
  \begin{center}\small
    This is the seal of Lecture 4: the server signs its part of the key exchange,\\
    so Mallory in the middle cannot replace it. Details of the protocol: Lecture 11.
  \end{center}
  \keymsg{On the web, a signature stops the man in the middle.}
  \say{TLS 1.3, protokół za kłódką w przeglądarce: serwer podpisuje cały przebieg uzgadniania klucza (wiadomość CertificateVerify, wykład 11). Bez tego wraca atak man in the middle z wykładu 4. Ale przeglądarka musi wiedzieć, że klucz publiczny naprawdę należy do tej strony. To jest pytanie na ostatnie slajdy.}
\end{frame}

\begin{frame}{Whose public key is it?}
  \vfill
  \begin{center}
    \begin{tikzpicture}[>=Stealth]
      \mallory{0,0}
      \node[hw, minimum width=5.8cm, minimum height=3.0cm] (srv) at (4.8,0) {};
      \node[font=\small\bfseries, text=nInk, anchor=north west] at ([xshift=4pt,yshift=-3pt]srv.north west) {\faServer\ update server};
      \node[msg, font=\scriptsize, minimum width=4.6cm, draw=sPink, line width=1.2pt] at (4.8,0.45) {\faFile*[regular]\ update.sh (Mallory's)};
      \node[msig, font=\scriptsize, minimum width=4.6cm] at (4.8,-0.2) {\faStamp\ update.sig (by Mallory)};
      \node[mpub, font=\scriptsize, minimum width=4.6cm] at (4.8,-0.85) {\faKey\ public\_key.pem (``Alice's'')};
      \draw[data] (M.east) -- (srv.west);
      \dave{10.8,0}
      \draw[net] (srv.east) -- (D.west);
      \node[font=\small, anchor=north] at (10.8,-1.3) {verify: \ok};
    \end{tikzpicture}
  \end{center}
  \vfill
  \begin{center}\small
    If Dave takes the public key from the same place as the update, Mallory replaces all three.\\
    Her signature fits her key, so \texttt{verify} says \texttt{True}.
  \end{center}
  \keymsg{A signature binds a message to a key; it cannot tell whose key it is.}
  \say{Ten sam problem co na początku ze skrótem, tylko o piętro wyżej. Przypomnieć punkt 2 slajdu What True means: podpis mówi o kluczu, nie o osobie. Jeśli klucz przyszedł z tego samego, zmienionego miejsca, to sprawdzenie jest bez wartości.}
\end{frame}

\begin{frame}{Task 9: whose key?}
  \vfill
  \taskgoal{as Dave's updater, check Alice's server and Mallory's server, first with the key from the server, then with the built-in key.}
  \taskready{\texttt{signing.py} (\texttt{verify}, \texttt{key\_hex}, \texttt{public\_key\_from\_hex}); two server folders with three files each.}
  \vfill
  \opentask{Task9}{about 7 minutes}
  \vfill
  \say{Około 7 minut. TODO 1 i 2: ten sam kod dla dwóch folderów, oba dają True. TODO 3: klucz z serwera Mallory to nie jest klucz wbudowany. TODO 4: z kluczem wbudowanym aktualizacja Mallory nie przechodzi. Zapytać na koniec: skąd program ma ten wbudowany klucz i dlaczego mu ufamy?}
\end{frame}

\samedot
\begin{frame}{Task 9: what we saw}
  % CHECKED: python3 Task9/solution/updater.py, 2026-11-17
  \vfill
  \begin{center}
    \renewcommand{\arraystretch}{1.45}
    \begin{tabular}{@{}l l c@{}}
      \toprule
      Alice's server & key from the server & \ok\ \texttt{True} \\
      Mallory's server & key from the server & \ok\ \texttt{True} \\
      the key from Mallory's server is Alice's key & & \nok\ \texttt{False} \\
      Mallory's server & key built into the app & \nok\ \texttt{False} \\
      \bottomrule
    \end{tabular}
  \end{center}
  \vfill
  \begin{center}\small The check is only as good as the way Dave got the public key.\end{center}
  \keymsg{We still need to know, safely, which public key is Alice's.}
  \say{Wynik rozwiązania. Z kluczem z serwera oba serwery przechodzą, więc sprawdzenie nic nie daje. Wbudowany klucz zatrzymuje Mallory. Ale wbudować klucz da się tylko, gdy znamy go z góry (aktualizacje jednej aplikacji). Przeglądarka łączy się z milionami stron: nie może mieć wbudowanych kluczy wszystkich.}
\end{frame}

\begin{frame}{A trusted party signs the key}
  \vfill
  \begin{center}
    \begin{tikzpicture}[>=Stealth]
      \node[hw, minimum width=5.6cm, minimum height=2.8cm] (c) at (0,0) {};
      \node[font=\small\bfseries, text=nInk, anchor=north] at (c.north) {``This key belongs to Alice''};
      \node[apub, font=\small] at (0,0) {\faKey\ Alice's public key};
      \node[hw, font=\scriptsize, fill=white] at (0,-0.85) {\faStamp\ signature of a trusted party};
      \node[font=\small, text=nInk, anchor=west, align=left] at (3.4,0) {Dave has the trusted party's key built in\\(a handful of keys, not millions),\\checks this signature, and then trusts\\Alice's key.};
    \end{tikzpicture}
  \end{center}
  \vfill
  \bigdef{A signed statement ``this public key belongs to this name'' is a \textbf{certificate}. Lecture 9.}
  \vfill
  \say{Most do wykładów 9 do 12 (PKI). Ten sam podpis z dzisiejszego wykładu, tylko podpisuje się klucz z nazwą. Zaufana strona to urząd certyfikacji (CA); jego klucz jest w systemie albo przeglądarce. Tu tylko idea; budowa certyfikatu X.509 i sprawdzanie na wykładzie 9.}
\end{frame}

\begin{frame}{Who wrote this message?}
  \vfill
  \begin{center}
    \renewcommand{\arraystretch}{1.6}
    \begin{tabular}{@{}l l l@{}}
      \toprule
      \textbf{what Bob gets} & \textbf{what it proves} & \textbf{to whom} \\
      \midrule
      a \digw{digest} & the file fits the digest; anybody could make both & nobody \\
      a MAC tag & made by someone with the shared \keyw{key} (Alice or Bob) & only Alice and Bob \\
      a \aliw{signature} & signed with \textbf{this private key}, unchanged since & \textbf{everybody} \\
      \bottomrule
    \end{tabular}
  \end{center}
  \vfill
  \bigdef{Who wrote it: the holder of the private key whose public key verifies the signature.\\
    That this key is Alice's we must know another way: a built-in key, or a certificate.}
  \vfill
  \say{Pełna odpowiedź na pytanie wykładu, z zastrzeżeniem z zadania 9. Tabela zestawia trzy narzędzia z dzisiejszego dnia: skrót (nic o autorze), MAC (jedna z dwóch osób, przekonuje tylko je), podpis (właściciel klucza, przekonuje każdego). Ostatnie zdanie to most do wykładu 9.}
\end{frame}

\begin{frame}{Good to know: quantum computers}
  \vfill
  \begin{columns}[c,onlytextwidth]
    \column{0.3\textwidth}
      \centering
      \begin{tikzpicture}
        \node[circle, fill=sPink, text=white, font=\fontsize{36}{38}\selectfont, minimum size=2.2cm] {\faAtom};
      \end{tikzpicture}
    \column{0.68\textwidth}
      \probbox{Shor's algorithm}{on a big quantum computer would break RSA, ECDSA and Ed25519: all three.}
      \stepbox{The new standard}{ML-DSA (FIPS 204, August 2024): a signature based on lattices, not on factoring or curves.}
      \plainbox{NIST plan (IR 8547): RSA, ECDSA and EdDSA deprecated after 2030, disallowed after 2035.}
  \end{columns}
  \vfill
  % CHECKED: FIPS 204 effective August 2024 (Federal Register 2024-17956); NIST IR 8547 (as in Lecture 5): ECDSA, EdDSA, RSA deprecated after 2030, disallowed after 2035
  \say{Można pominąć. Dla porównania: klucz publiczny ML-DSA-44 ma 1312 bajtów, podpis 2420 bajtów, czyli dużo więcej niż Ed25519, ale interfejs jest ten sam: sign i verify. Takiego komputera kwantowego jeszcze nie ma. Plan NIST jest ten sam co w wykładach 3 i 5.}
\end{frame}

\end{document}
